\documentclass[11pt]{article}
% Shared preamble for the rigor documents (Lamport-style structured proofs).  \input this file after \documentclass.
\usepackage[T1]{fontenc}\usepackage{lmodern}\usepackage{microtype}
\usepackage[margin=1in]{geometry}
\usepackage{amsmath,amssymb,amsthm,mathtools}
\usepackage{xcolor}\usepackage{enumitem}\usepackage{booktabs}\usepackage{graphicx}\usepackage{hyperref}
\usepackage[most]{tcolorbox}
\definecolor{navy}{HTML}{17365D}\definecolor{teal}{HTML}{117A8B}\definecolor{warn}{HTML}{A33A2B}\definecolor{okgreen}{HTML}{2E7D4F}
\hypersetup{colorlinks=true,linkcolor=navy,citecolor=teal,urlcolor=teal}
\theoremstyle{plain}
\newtheorem{theorem}{Theorem}[section]\newtheorem{lemma}[theorem]{Lemma}\newtheorem{proposition}[theorem]{Proposition}\newtheorem{corollary}[theorem]{Corollary}
\theoremstyle{definition}\newtheorem{definition}[theorem]{Definition}\newtheorem{remark}[theorem]{Remark}\newtheorem{claim}[theorem]{Claim}\newtheorem{issue}[theorem]{Issue}
% ---------- Lamport structured proofs ----------
% \lstep{level}{number}{assertion}   -> indented "<level>number. assertion"
% \lproof                             -> "PROOF:" line (start of the sub-proof of the preceding step)
% \lby{justification}                 -> "BY ..." leaf justification for the preceding step
% \lqed{level}{number}                -> QED step of a sub-proof
% keywords: \lassume \lprove \lsuffices \lcase \ldefine \lpick \lnew
\newlength{\lindent}\setlength{\lindent}{1.7em}
\newcommand{\lstep}[3]{\par\smallskip\noindent\hspace*{\dimexpr\numexpr#1-1\relax\lindent\relax}{\color{navy}$\langle#1\rangle#2$.}\ #3}
\newcommand{\lproof}{\par\noindent\hspace*{\lindent}{\scshape Proof:}}
\newcommand{\lby}[1]{\par\noindent\hspace*{\lindent}{\scshape By}\ #1.}
\newcommand{\lqed}[2]{\par\smallskip\noindent\hspace*{\dimexpr\numexpr#1-1\relax\lindent\relax}{\color{navy}$\langle#1\rangle#2$.}\ {\scshape Q.E.D.}}
\newcommand{\lassume}{{\scshape Assume}\ }\newcommand{\lprove}{{\scshape Prove}\ }\newcommand{\lsuffices}{{\scshape Suffices}\ }
\newcommand{\lcase}{{\scshape Case}\ }\newcommand{\ldefine}{{\scshape Define}\ }\newcommand{\lpick}{{\scshape Pick}\ }\newcommand{\lnew}{{\scshape New}\ }
% ---------- verdict boxes ----------
\newtcolorbox{verdictok}{colback=okgreen!8,colframe=okgreen,title={\bfseries Verdict: verified},fonttitle=\small}
\newtcolorbox{verdictgap}{colback=orange!10,colframe=orange!80!black,title={\bfseries Verdict: gap / caveat},fonttitle=\small}
\newtcolorbox{verdictbreak}{colback=warn!10,colframe=warn,title={\bfseries Verdict: proof-breaking issue},fonttitle=\small}
\newtcolorbox{citedbox}[1]{colback=navy!5,colframe=navy!60,title={\bfseries Cited result: #1},fonttitle=\small,breakable}
\setlength{\parindent}{0pt}\setlength{\parskip}{4pt}\allowdisplaybreaks
\newcommand{\cA}{\mathcal A}\newcommand{\cF}{\mathcal F}\newcommand{\cM}{\mathfrak M}\newcommand{\cN}{\mathfrak N}

% extra calligraphic shortcuts (\providecommand: no clash if a document defines its own)
\providecommand{\cB}{\mathcal B}\providecommand{\cC}{\mathcal C}\providecommand{\cD}{\mathcal D}\providecommand{\cE}{\mathcal E}\providecommand{\cG}{\mathcal G}\providecommand{\cH}{\mathcal H}\providecommand{\cI}{\mathcal I}\providecommand{\cJ}{\mathcal J}\providecommand{\cK}{\mathcal K}\providecommand{\cL}{\mathcal L}\providecommand{\cO}{\mathcal O}\providecommand{\cP}{\mathcal P}\providecommand{\cQ}{\mathcal Q}\providecommand{\cR}{\mathcal R}\providecommand{\cS}{\mathcal S}\providecommand{\cT}{\mathcal T}\providecommand{\cU}{\mathcal U}\providecommand{\cV}{\mathcal V}\providecommand{\cW}{\mathcal W}\providecommand{\cX}{\mathcal X}\providecommand{\cY}{\mathcal Y}\providecommand{\cZ}{\mathcal Z}
\newcommand{\Fid}{\operatorname{F}}\newcommand{\Tr}{\operatorname{Tr}}\newcommand{\eps}{\varepsilon}\newcommand{\dd}{\,\mathrm d}

\newcommand{\hh}{\mathfrak h}
\newcommand{\Erecc}{E_{\mathrm{rec}}}
\newcommand{\Stwo}{\mathfrak S_2}
\newcommand{\Herm}{\operatorname{Herm}}
\newcommand{\SLD}{\mathcal S_Q}
\newcommand{\Aop}{\mathbb A}

\title{The exact second-order optimum over \emph{all} quasi-free channels:\\
the tangent problem, its KKT sign criterion, and what is left open}
\author{Track F1 (Phase 5, analytic)}
\date{2026-09-15}

\begin{document}
\maketitle

\begin{abstract}
The second-order recovery problem for the chiral free fermion is the convex programme
$E^{(2)}=\min_{(X,Z)\in\cF}g_Q(Q-Q_{\rm rec}(X,Z))$ (S8 Thm.~3.1, S9 Thm.~3.5).  Its value
on the \emph{isometric orbit} is $(1-\theta)f_2z^2$ (S10 Lemma~6.1, Thm.~6.2); whether the
non-isometric data (loss of isometry $N$, noise $Y$) can do better was left open (S10,
final Remark).  We give (i) an exact parametrisation of $\cF$ by the polar co-isometry $V$
(which exists iff $1\notin\sigma_p(N)$), a positive $N$ and a noise $Y$ with $0\le Y\le N$;
(ii) the tangent (leading-order) convex programme in the variables $(\zeta,N_1,Y_1)$, with
every block and every sign explicit, and a constructive proof that its value is an upper
bound for $\limsup z^{-2}\min_\cF g_Q$; (iii) the \emph{inclusion} of the boundary-moving
(non-surjective) isometry directions in the cone generated by the anti-symmetric parts of admissible isometry
generators and the singular $N_1\ge0$ directions --- an inclusion which cannot be reversed
(reversing it would give $T=0$); (iv) the KKT sign criterion at the isometric optimum, in
the exact form $\Tr(N_1M_*)\ge\Tr[(N_1^{1/2}\tau_*N_1^{1/2})_+]$ for all $N_1\ge0$, with
$M_*=E_D\Herm(t_*Q)E_D$ and $\tau_*=E_Dt_*E_D$ read as quadratic forms, its equivalence to
the pair $M_*\ge0$, $M_*\ge\tau_*$ (condition (a)), and the \emph{endpoint} condition (e)
that must be added to make it a complete optimality criterion; (v) the identification of
these with S10's multipliers $\Sigma_2,\Sigma_1$; (vi) a reduction of $\theta$ to a
geometry-free Wiener--Hopf problem.  The lower bound (interchange of $\min$ and $z\to0$) is
isolated as Hypothesis~H2, and the reduction of (e) to the two model directions as
Hypothesis~H3.  This version incorporates the repairs R1--R15 of the referee
report \texttt{referee\_exact\_optimum\_tangent\_problem.tex} and the repairs S1--S6 of
\texttt{referee2\_exact\_optimum\_tangent\_problem.tex} (both 2026-09-15), together with
the numerical inputs of tracks F2 and F3.
\end{abstract}

\tableofcontents

\section{Setting, notation, and the two conventions that must be fixed}\label{sec:setup}
\paragraph{(C1) Geometry and one-particle space.}
We use the half-line normalisation of S10~(N5), in which the collar parameter and the
cross ratio coincide.  $H:=L^2(\mathbb R)$, $P$ the chiral vacuum (Hardy) projection,
$P^\perp:=1-P$; $E_S$ is multiplication by $\mathbf 1_S$.  Fixed:
$A=(-\infty,0)$, $D=(0,1)$, $I=A\cup D=(-\infty,1)$.  Moving with the parameter $s>0$:
$B=(0,\tfrac1{1+s})$, $C=(\tfrac1{1+s},1)$, $D=B\cup C$.  We write
$\hh_S:=L^2(S)$, so $\hh_I=\hh_A\oplus\hh_D$, $\hh_D=\hh_B\oplus\hh_C$, and all
\emph{symbols} act on $\hh_I$.  The vacuum symbol is $Q=E_IPE_I$, $0<Q<1$ (S10~(C2)), with
blocks $Q_{AA},Q_{AD},Q_{DA},Q_{DD}$ and $Q_{BB}=E_BPE_B$; $R_\pm$ denotes $Q_{BB}$ resp.\
$1_{\hh_B}-Q_{BB}$.  For the chiral free fermion $z=s(1+o(1))$ and $f_2=c/(12\pi^2)$.

\paragraph{(C2) Feasible set and recovered symbol.}  With $X:\hh_B\to\hh_D$ bounded and
$Z=Z^*$ on $\hh_D$,
\begin{equation}\label{eq:F}
 \cF:=\{(X,Z):\ XR_+X^*\le Z\le 1_{\hh_D}-XR_-X^*\},\qquad
 Q_{\rm rec}(X,Z)=\begin{pmatrix}Q_{AA}&Q_{AB}X^*\\ XQ_{BA}&Z\end{pmatrix},
\end{equation}
(S8 Thm.~3.1, S9 Thm.~3.5: \emph{every} recovery channel's recovered symbol lies in $\cF$,
and every point of $\cF$ is realised).  The defect is $\delta:=Q-Q_{\rm rec}(X,Z)$; note
$\delta_{AA}=0$ identically.

\paragraph{(C3) Notation dictionary.}  S10 writes $V:=X^*:\hh_D\to\hh_B$ for the Heisenberg
one-particle map.  In this note $V$ denotes the \emph{co-isometric factor} of $X$ in
Lemma~\ref{lem:polar}, so that S10's $V$ equals our $V^*(1-N)^{1/2}=X^*$.  The compression
is $W:=E_BW_sE_D:\hh_D\to\hh_B$ (S10's $V_c$), a \emph{unitary} $\hh_D\to\hh_B$ (S10~(C5)),
where $W_s$ is the half-density push-forward along $x\mapsto x/(1+sx)$ on $D$ and the
identity on $A$, and
\begin{equation}\label{eq:Wexp}
 W_s=1-sD_\chi+O(s^2),\qquad D_\chi=\chi\partial_x+\tfrac12\chi',\quad
 \chi(x)=-x^2\mathbf 1_{x>0}\in C^{1,1},\quad D_\chi^*=-D_\chi .
\end{equation}
The compression point of $\cF$ is $(X_c,Z_c)=(W^*,\,W^*Q_{BB}W)$ and its defect is
$\delta_c=E_I(P-W_s^*PW_s)E_I=s\dot\delta+O(s^2)$ with
$\dot\delta:=E_I\Theta E_I$, $\Theta:=[P,D_\chi]$ (S10~(N5)).

\paragraph{(C4) Objective, SLD, and the normalisation of the gradient.}
In an eigenbasis of $Q$ (eigenvalues $q_i$, modular variables $q=(1+e^\kappa)^{-1}$),
$w_{ik}:=[q_i(1-q_k)+q_k(1-q_i)]^{-1}\ge1$ and
\begin{equation}\label{eq:gQ}
 g_Q(\delta)=\tfrac14\sum_{ik}w_{ik}|\widetilde\delta_{ik}|^2
 =\tfrac18\sup_{t=t^*}\frac{(\Tr t\delta)^2}{V(t)},\qquad V(t)=\Tr\bigl(tQt(1-Q)\bigr)
 =\|P^\perp tP\|_2^2 ,
\end{equation}
the supremum being attained at the symmetric logarithmic derivative (SLD)
$t_\delta:=\SLD^{-1}\delta$, $\SLD(t):=Qt(1-Q)+(1-Q)tQ$ (S8 eq.~(3.9), S10~(C4),(N4)).
$g_Q$ is a positive quadratic form on the real vector space of self-adjoint symbols; we
write $b$ for its polarisation, $g_Q(\delta)=b(\delta,\delta)$, and
\begin{equation}\label{eq:formdom}
 \mathfrak F:=\{\eta=\eta^*=E_I\eta E_I:\ g_Q(\eta)<\infty\}
\end{equation}
for its form domain, a real Hilbert space.  \emph{All} statements below are inside
$\mathfrak F$.

\begin{lemma}[Normalisation]\label{lem:norm}
Let $\delta\in\mathfrak F$ and $t_\delta=\SLD^{-1}\delta$.  Then, with $\eta\in\mathfrak F$,
\begin{equation}\label{eq:normal}
 \widetilde{(t_\delta)}_{ik}=w_{ik}\widetilde\delta_{ik},\qquad
 g_Q(\delta)=\tfrac14\Tr(t_\delta\delta),\qquad V(t_\delta)=2g_Q(\delta),\qquad
 b(\delta,\eta)=\tfrac14\Tr(t_\delta\eta),
\end{equation}
hence $\frac{\dd}{\dd \sigma}g_Q(\delta+\sigma\eta)\big|_{\sigma=0}=\tfrac12\Tr(t_\delta\eta)$
and $g_Q(\delta)=\frac18(\Tr t_\delta\delta)^2/V(t_\delta)$, consistent with \eqref{eq:gQ}.
\end{lemma}
\begin{proof}
\lstep{1}{1}{$(\SLD t)^{\sim}_{ik}=\widetilde t_{ik}/w_{ik}$, hence the first identity.}
\lby{$\SLD(t)=Qt(1-Q)+(1-Q)tQ$ acts diagonally in the eigenbasis of $Q$ with eigenvalue
 $q_i(1-q_k)+q_k(1-q_i)=w_{ik}^{-1}$}
\lstep{1}{2}{$\Tr(t_\delta\eta)=\sum_{ik}w_{ik}\widetilde\delta_{ik}\overline{\widetilde\eta_{ik}}=4b(\delta,\eta)$
 for self-adjoint $\delta,\eta$; in particular $\Tr(t_\delta\delta)=4g_Q(\delta)$.}
\lby{$\langle1\rangle1$, $\Tr(t\eta)=\sum_{ik}\widetilde t_{ik}\widetilde\eta_{ki}$,
 $\widetilde\eta_{ki}=\overline{\widetilde\eta_{ik}}$, and \eqref{eq:gQ}}
\lstep{1}{3}{$V(t_\delta)=\sum_{ik}|\widetilde{(t_\delta)}_{ik}|^2q_k(1-q_i)
 =\tfrac12\sum_{ik}|\widetilde{(t_\delta)}_{ik}|^2w_{ik}^{-1}=\tfrac12\sum_{ik}w_{ik}|\widetilde\delta_{ik}|^2=2g_Q(\delta)$.}
\lby{\eqref{eq:gQ}, symmetrising $q_k(1-q_i)$ in $i\leftrightarrow k$, and $\langle1\rangle1$}
\lstep{1}{4}{The derivative formula and the Legendre check.}
\lby{$g_Q(\delta+\sigma\eta)=g_Q(\delta)+2\sigma b(\delta,\eta)+\sigma^2g_Q(\eta)$ with
 $\langle1\rangle2$; and $(4g)^2/(8\cdot2g)=g$}
\lqed{1}{5}
\end{proof}

\begin{remark}[Factor $2$ in S10]\label{rem:factor2}
S10~(C4) states $\Tr(t_\delta\delta)=2g_Q(\delta)$, $V(t_\delta)=\frac12g_Q(\delta)$, which
corresponds to the convention $\SLD(t)=\frac12\delta$ (equivalently $2\SLD(t)=\delta$,
i.e.\ $t\mapsto t/2$).  \emph{Erratum (referee, 2026-09-15)}: S10~(C4) is internally
inconsistent, since it also states $\SLD(t)=\delta$ (its eq.~(5)); with that equation the
correct relations are the ones in \eqref{eq:normal}.  No S10 conclusion is affected --- the
Legendre ratio \eqref{eq:gQ} and every sign condition below are homogeneous of degree $1$
in $t_*$ --- but any numerical work must say which normalisation its ``units of
$g_Q(\dot\delta)$'' refer to.  We use \eqref{eq:normal} throughout.
\end{remark}


\section{Exact parametrisation of the feasible set}\label{sec:param}
Throughout this section $s>0$ is fixed; nothing is expanded.

\begin{lemma}[The $(N,Y)$ chart]\label{lem:NY}
Let $X:\hh_B\to\hh_D$ be bounded and $Z=Z^*$ on $\hh_D$.  Put
\begin{equation}\label{eq:NY}
 N:=1_{\hh_D}-XX^* ,\qquad Y:=Z-XQ_{BB}X^* .
\end{equation}
Then $(X,Z)\in\cF$ if and only if $N\ge0$ and $0\le Y\le N$.  The map
$(X,Z)\mapsto(X,Y)$ is an affine bijection of $\cF$ onto
$\{(X,Y):\ 0\le Y\le 1-XX^*\}$.
\end{lemma}
\begin{proof}
\lstep{1}{1}{$Z\ge XR_+X^*\iff Y\ge0$.}
\lby{definition \eqref{eq:NY} of $Y$ with $R_+=Q_{BB}$}
\lstep{1}{2}{$Z\le1-XR_-X^*\iff Y\le N$.}
\lby{$1-XR_-X^*-Z=1-XR_-X^*-XR_+X^*-Y=1-XX^*-Y=N-Y$, using $R_++R_-=1_{\hh_B}$}
\lstep{1}{3}{If $(X,Z)\in\cF$ then $N\ge0$; conversely $N\ge0$ and $0\le Y\le N$ give
 $(X,Z)\in\cF$.}
\lby{$\langle1\rangle1$--$\langle1\rangle2$: $0\le Y\le N$ forces $N\ge Y\ge0$; the
 converse is the conjunction of the two displayed equivalences}
\lstep{1}{4}{Bijectivity and affineness.}
\lby{$Z=XQ_{BB}X^*+Y$ for fixed $X$, an affine bijection in the second variable}
\lqed{1}{5}
\end{proof}

\begin{lemma}[Polar decomposition; the co-isometric factor]\label{lem:polar}
Let $X:\hh_B\to\hh_D$ with $N=1-XX^*\ge0$.  Then $X=(1-N)^{1/2}V$ where $V:\hh_B\to\hh_D$
is the partial isometry of the polar decomposition of $X$, with initial space
$(\ker X)^\perp$ and final space $\overline{\operatorname{ran}X}=\overline{\operatorname{ran}(1-N)}$.  Moreover
\begin{enumerate}[leftmargin=1.6em,itemsep=1pt]
\item $VV^*=1_{\hh_D}-\mathbf 1_{\{1\}}(N)$; hence $V$ is a co-isometry ($V^*$ an isometry
 $\hh_D\to\hh_B$) if and only if $1\notin\sigma_p(N)$, i.e.\ iff $X$ has dense range;
\item if in addition $\ker X=0$, $V$ is unitary;
\item in particular, if $\|N\|<1$ then $V$ is a co-isometry, and $X\mapsto(N,V)$ is a
 bijection from $\{X:\|1-XX^*\|<1\}$ onto $\{(N,V):0\le N,\ \|N\|<1,\ VV^*=1\}$ modulo the
 free choice of $V$ on $\ker X$ (empty when $\ker X=0$);
\item statements (1)--(3) concern the \emph{polar} $V$.  The existence of \emph{some}
 co-isometric factor, i.e.\ of a co-isometry $V'$ with $X=(1-N)^{1/2}V'$, is the weaker
 condition $\dim\ker(1-N)\le\dim\ker X$ (extend the polar $V$ from $(\ker X)^\perp$ by any
 isometry $\ker(1-N)\to\ker X$; the extension changes $X$ by $0$ because
 $(1-N)^{1/2}$ annihilates $\ker(1-N)$).  Only the polar $V$ is used below.
\end{enumerate}
\end{lemma}
\begin{proof}
\lstep{1}{1}{$X=U|X|$ with $|X|=(X^*X)^{1/2}$ and $U$ a partial isometry with initial space
 $(\ker X)^\perp$, final space $\overline{\operatorname{ran}X}$; and $X=|X^*|U$ with
 $|X^*|=(XX^*)^{1/2}=(1-N)^{1/2}$.  Set $V:=U$.}
\lby{polar decomposition; $|X^*|U=U|X|U^*U=U|X|=X$ because $U|X|U^*=|X^*|$}
\lstep{1}{2}{$VV^*=$ the orthogonal projection onto $\overline{\operatorname{ran}X}
 =\overline{\operatorname{ran}XX^*}=\overline{\operatorname{ran}(1-N)}=(\ker(1-N))^\perp$,
 and $\ker(1-N)=\mathbf 1_{\{1\}}(N)\hh_D$.}
\lby{for a partial isometry, $VV^*$ is the projection onto the final space;
 $\overline{\operatorname{ran}X}=\overline{\operatorname{ran}XX^*}$ for any bounded $X$; the spectral theorem for $N\ge0$}
\lstep{1}{3}{(1) and (2).}
\lby{$\langle1\rangle2$: $VV^*=1$ iff $\mathbf 1_{\{1\}}(N)=0$; $V^*V$ is the projection onto
 $(\ker X)^\perp$, so $V$ is unitary iff also $\ker X=0$}
\lstep{1}{4}{(3).}
\lby{$\|N\|<1\Rightarrow1\notin\sigma(N)$; conversely $(N,V)$ with the stated properties
 gives $X:=(1-N)^{1/2}V$ with $XX^*=(1-N)^{1/2}VV^*(1-N)^{1/2}=1-N$}
\lqed{1}{5}
\end{proof}

\begin{remark}[The co-isometry is not automatic]\label{rem:shift}
Statement (1) is sharp: for $\hh_B=\hh_D=\ell^2(\mathbb N)$ and $X=S$ the forward shift,
$N=1-SS^*$ is the rank-one projection onto $e_0$, $1\in\sigma_p(N)$, and no co-isometry $V$
with $X=(1-N)^{1/2}V$ exists (such a $V$ would satisfy $V=S+e_0\varphi^*$ and then
$VV^*=1$ forces both $S\varphi=0$ and $\|\varphi\|=1$).  This is why the tangent analysis
of \S\ref{sec:tangent} is carried out in the regime $\|N\|<1$, where Lemma~\ref{lem:polar}(3)
applies; it covers a norm neighbourhood of the compression, where $\ker X=0$ as well.
\end{remark}

\begin{lemma}[Factorisation through the compression]\label{lem:WU}
Let $V:\hh_B\to\hh_D$ be a co-isometry and $W:\hh_D\to\hh_B$ the compression unitary
of~(C3).  Then $U:=W^*V^*:\hh_D\to\hh_D$ is an isometry, $V^*=WU$, and $U$ is onto iff
$V^*$ is onto (iff $V$ is unitary).  Conversely $V^*:=WU$ is a co-isometry's adjoint for
every isometry $U$ of $\hh_D$.
\end{lemma}
\begin{proof}
\lstep{1}{1}{$U^*U=VWW^*V^*=VV^*=1_{\hh_D}$, and $WU=WW^*V^*=V^*$.}
\lby{$W$ unitary $\hh_D\to\hh_B$ ($WW^*=1_{\hh_B}$, $W^*W=1_{\hh_D}$) and $VV^*=1_{\hh_D}$}
\lstep{1}{2}{$\operatorname{ran}U=W^*\operatorname{ran}V^*$ and $W^*$ is unitary; the converse.}
\lby{$\langle1\rangle1$ and unitarity of $W^*$ for the first claim; for the converse put
 $V:=U^*W^*:\hh_B\to\hh_D$, then $V^*=WU$ and $VV^*=U^*W^*WU=U^*U=1_{\hh_D}$}
\lqed{1}{3}
\end{proof}



\begin{proposition}[The defect in the $(V,N,Y)$ chart, exactly]\label{prop:blocks}
Let $(X,Z)\in\cF$ with $\|N\|<1$, $X=(1-N)^{1/2}V$ as in Lemma~\ref{lem:polar}(3), and
$Y=Z-XQ_{BB}X^*\in[0,N]$.  Then $\delta=Q-Q_{\rm rec}(X,Z)$ has the blocks
\begin{equation}\label{eq:exactblocks}
\begin{aligned}
 \delta_{AA}&=0,\\
 \delta_{AD}&=Q_{AD}-Q_{AB}V^*(1-N)^{1/2}
  =Q_{AD}-\Lambda\,U(1-N)^{1/2},\qquad \Lambda:=Q_{AB}W:\hh_D\to\hh_A,\\
 \delta_{DA}&=\delta_{AD}^*,\\
 \delta_{DD}&=Q_{DD}-(1-N)^{1/2}V\,Q_{BB}\,V^*(1-N)^{1/2}-Y
  =Q_{DD}-(1-N)^{1/2}U^*\!\bigl(W^*Q_{BB}W\bigr)U(1-N)^{1/2}-Y ,
\end{aligned}
\end{equation}
with $V^*=WU$ as in Lemma~\ref{lem:WU}.  In particular $W^*Q_{BB}W=Z_c$ is the
compression's $DD$-block, and at $(N,Y,U)=(0,0,1)$ one recovers $\delta=\delta_c$.
\end{proposition}
\begin{proof}
\lstep{1}{1}{$\delta_{AA}=Q_{AA}-Q_{AA}=0$ and $\delta_{AD}=Q_{AD}-Q_{AB}X^*$.}
\lby{\eqref{eq:F}}
\lstep{1}{2}{$X^*=V^*(1-N)^{1/2}=WU(1-N)^{1/2}$, hence the two forms of $\delta_{AD}$.}
\lby{Lemma~\ref{lem:polar} ($(1-N)^{1/2}$ is self-adjoint) and Lemma~\ref{lem:WU}}
\lstep{1}{3}{$\delta_{DD}=Q_{DD}-Z=Q_{DD}-XQ_{BB}X^*-Y$ and
 $XQ_{BB}X^*=(1-N)^{1/2}VQ_{BB}V^*(1-N)^{1/2}=(1-N)^{1/2}U^*W^*Q_{BB}WU(1-N)^{1/2}$.}
\lby{Lemma~\ref{lem:NY} and $\langle1\rangle2$, using $V=U^*W^*$}
\lstep{1}{4}{The compression point.}
\lby{$(N,Y,U)=(0,0,1)$ gives $X=V=W^*$, i.e.\ $(X,Z)=(X_c,Z_c)$ of~(C3)}
\lqed{1}{5}
\end{proof}

\begin{remark}[Reading \eqref{eq:exactblocks}]\label{rem:read}
$Q_{AD}$ and $Q_{AB}$ are different operators ($Q_{AD}=E_APE_D$, $Q_{AB}=E_APE_B$), so
the abbreviation ``$\delta_{AD}=Q_{AB}(1-V^*(1-N)^{1/2})$'' used in the Phase-5 brief is to
be read through $\Lambda=Q_{AB}W$: the first form in \eqref{eq:exactblocks} is the exact
one.  As $s\downarrow0$, $B\uparrow D$, $W\to1$ and $\Lambda\to Q_{AD}$; this is the only
place where the geometry enters the algebra.
\end{remark}

\begin{verdictok}[\;\S\ref{sec:param}]
Lemmas~\ref{lem:NY}--\ref{lem:WU} and Proposition~\ref{prop:blocks} are exact: no
expansion, no smallness, no dimension assumption beyond
$\dim\hh_B=\dim\hh_D=\aleph_0$ (used only through Lemma~\ref{lem:polar}(3), i.e.\ through
$\|N\|<1$).  The chart is
$\cF\cap\{\|N\|<1\}\ \cong\ \{(U\ \text{isometry of }\hh_D,\ 0\le N<1,\ 0\le Y\le N)\}$.
\end{verdictok}

\section{The tangent problem}\label{sec:tangent}
\subsection{The scaling regime and the first-order defect}\label{sec:scaling}

\begin{definition}[Admissible first-order data]\label{def:data}
A \emph{first-order datum} is a triple $(\zeta,N_1,Y_1)$ where
\begin{enumerate}[leftmargin=1.6em,itemsep=1pt]
\item $\zeta$ is the generator of a $C_0$-\emph{semigroup} $(U_s)_{s\ge0}$ of
 \emph{isometries} of $\hh_D$ with $U_0=1$, $U_s=1+s\zeta+o(s)$ strongly on
 $\operatorname{dom}\zeta$ (so $U_s^*U_s=1_{\hh_D}$; $U_s$ need not be onto).  The
 semigroup property is used only through Lumer--Phillips in Thm.~\ref{thm:bdry}(1); it is
 satisfied by every family occurring here ($e^{s\mathcal G}$ and the push-forwards
 $W_{\varphi_s}$ along a flow).  Equivalently one may assume that $s\mapsto U_s^*g$ is
 weakly differentiable;
\item $N_1$ is a \emph{bounded} operator with $N_1\ge0$, and $0\le Y_1\le N_1$
 (boundedness is needed for $(1-sN_1)^{1/2}$ in \eqref{eq:scaled} and for
 Lemma~\ref{lem:feas}; positive \emph{forms} $N_1$ occur only as limits, in
 Thm.~\ref{thm:bdry} and in condition~(e) of Thm.~\ref{thm:kktmain}, where no scaled
 family is constructed);
\item the operator $\delta^{(1)}$ of \eqref{eq:delta1} lies in the form domain
 $\mathfrak F$ of \eqref{eq:formdom}.
\end{enumerate}
The associated \emph{scaled family} is, for $0<s<\|N_1\|^{-1}$,
\begin{equation}\label{eq:scaled}
 N(s):=sN_1,\qquad Y(s):=sY_1,\qquad V(s)^*:=W_sU_s,\qquad
 X(s):=(1-sN_1)^{1/2}V(s),\qquad Z(s):=X(s)Q_{BB}X(s)^*+sY_1 .
\end{equation}
\end{definition}

\begin{lemma}[The scaled family is exactly feasible]\label{lem:feas}
For every first-order datum and every $0<s<\|N_1\|^{-1}$, $(X(s),Z(s))\in\cF$, with
$1-X(s)X(s)^*=sN_1$ and $Y=sY_1$.
\end{lemma}
\begin{proof}
\lstep{1}{1}{$V(s)V(s)^*=U_s^*W_s^*W_sU_s=U_s^*U_s=1_{\hh_D}$, so $V(s)$ is a co-isometry.}
\lby{$W_s$ restricted as in~(C3) is unitary $\hh_D\to\hh_B$ and $U_s$ is an isometry}
\lstep{1}{2}{$X(s)X(s)^*=(1-sN_1)^{1/2}V(s)V(s)^*(1-sN_1)^{1/2}=1-sN_1$, hence $N=sN_1\ge0$.}
\lby{$\langle1\rangle1$; $1-sN_1\ge0$ because $s\|N_1\|<1$}
\lstep{1}{3}{$0\le Y=sY_1\le sN_1=N$, hence $(X(s),Z(s))\in\cF$.}
\lby{$0\le Y_1\le N_1$ and $s>0$; Lemma~\ref{lem:NY}}
\lqed{1}{4}
\end{proof}

\begin{lemma}[The $AA$- and $DD$-blocks of $\dot\delta$ vanish]\label{lem:diagzero}
$E_A\dot\delta E_A=0$ and $E_D\dot\delta E_D=0$.
\end{lemma}
\begin{proof}
\lstep{1}{1}{$D_\chi E_A=0=E_AD_\chi$.}
\lby{$\chi=-x^2\mathbf 1_{x>0}$ vanishes identically on $\overline A=(-\infty,0]$ together
 with $\chi'$, and $D_\chi=\chi\partial_x+\frac12\chi'$ is a local (first-order
 differential) operator, so $D_\chi f=0$ for $f$ supported in $\overline A$ and no boundary
 distribution arises at $0$ since $\chi(0)=\chi'(0)=0$; the second identity is the adjoint
 of the first with $D_\chi^*=-D_\chi$}
\lstep{1}{2}{$E_A\dot\delta E_A=E_A(PD_\chi-D_\chi P)E_A=(E_AP)(D_\chi E_A)-(E_AD_\chi)(PE_A)=0$.}
\lby{$\langle1\rangle1$ and $\dot\delta=E_I\Theta E_I$, $\Theta=[P,D_\chi]$, $E_AE_I=E_A$}
\lstep{1}{3}{$E_D\dot\delta E_D=0$.}
\lby{S10 Lemma~11.1 (M\"obius rigidity): $\chi|_D$ coincides with the global M\"obius field
 $-x^2$, whose generator commutes with $P$, so $[p,d]=0$ as a form on $C_c^\infty(D)$, and
 $\dot\delta$ is bounded}
\lqed{1}{4}
\end{proof}

\begin{proposition}[First-order defect: all blocks, all signs]\label{prop:delta1}
Let $(\zeta,N_1,Y_1)$ be a first-order datum and $(X(s),Z(s))$ its scaled family.  Then, as
sesquilinear forms on $C_c^\infty(A)\oplus C_c^\infty(D)$,
\begin{equation}\label{eq:delta1}
 \delta(s)=s\,\delta^{(1)}+o(s),\qquad
 \delta^{(1)}=\dot\delta+\Aop(\zeta,N_1,Y_1),
\end{equation}
where $\Aop$ is the real-linear map with blocks
\begin{equation}\label{eq:Ablocks}
\begin{aligned}
 \Aop_{AA}&=0, &
 \Aop_{AD}&=-\,Q_{AD}\,\zeta\;+\;\tfrac12\,Q_{AD}N_1 , &
 \Aop_{DA}&=\Aop_{AD}^*=-\,\zeta^*Q_{DA}+\tfrac12N_1Q_{DA},\\[2pt]
 \Aop_{DD}&=-\bigl(\zeta^*Q_{DD}+Q_{DD}\zeta\bigr)\;+\;\tfrac12\{N_1,Q_{DD}\}\;-\;Y_1 . &&&&
\end{aligned}
\end{equation}
The signs are: the noise $Y_1$ enters $\delta_{DD}$ with a \emph{minus}; the loss of
isometry $N_1$ enters $\delta_{DD}$ with $+\frac12\{N_1,Q_{DD}\}$ and $\delta_{DA}$ with
$+\frac12N_1Q_{DA}$.  Moreover $\dot\delta$ has only the blocks $\dot\delta_{AD}$,
$\dot\delta_{DA}$ (Lemma~\ref{lem:diagzero}).  The statement is one of sesquilinear forms
only; what upgrades \eqref{eq:delta1} to a statement in the $g_Q$-norm --- which is
strictly stronger, the weights $w_{ik}$ being unbounded --- is Hypothesis~H1
(Def.~\ref{hyp:H1}).
\end{proposition}
\begin{proof}
\lstep{1}{1}{$\delta_{AD}(s)=(Q_{AD}-\Lambda_s)+\Lambda_s\bigl(1-U_s(1-sN_1)^{1/2}\bigr)$ with
 $\Lambda_s=Q_{AB}W_s$.}
\lby{Proposition~\ref{prop:blocks} and adding and subtracting $\Lambda_s$}
\lstep{1}{2}{$Q_{AD}-\Lambda_s=s\dot\delta_{AD}+O(s^2)$, $\Lambda_s=Q_{AD}+O(s)$, and
 $1-U_s(1-sN_1)^{1/2}=-s\zeta+\frac s2N_1+o(s)$.}
\lby{$\Lambda_s$ is the $AD$-block of $Q_{\rm rec}(X_c,Z_c)$ and
 $\delta_c=s\dot\delta+O(s^2)$ (C3); $U_s=1+s\zeta+o(s)$ and
 $(1-sN_1)^{1/2}=1-\frac s2N_1+O(s^2)$ by the spectral theorem}
\lstep{1}{3}{$\delta_{AD}(s)=s\bigl[\dot\delta_{AD}-Q_{AD}\zeta+\frac12Q_{AD}N_1\bigr]+o(s)$.}
\lby{$\langle1\rangle1$--$\langle1\rangle2$: the $O(s)$ error in $\Lambda_s$ multiplies a
 factor $O(s)$}
\lstep{1}{4}{$\delta_{DD}(s)=(Q_{DD}-Z_c)-s(\zeta^*Z_c+Z_c\zeta)+\frac s2\{N_1,Z_c\}-sY_1+o(s)$.}
\lby{Proposition~\ref{prop:blocks} with $W^*Q_{BB}W=Z_c$, expanding
 $(1-\frac s2N_1)(1+s\zeta^*)Z_c(1+s\zeta)(1-\frac s2N_1)$ and keeping first order}
\lstep{1}{5}{$Q_{DD}-Z_c=O(s^2)=o(s)$ and $Z_c=Q_{DD}+o(1)$, hence \eqref{eq:Ablocks}.}
\lby{$Q_{DD}-Z_c=(\delta_c)_{DD}=s\,E_D\dot\delta E_D+O(s^2)=O(s^2)$ by
 Lemma~\ref{lem:diagzero}; then $\langle1\rangle4$}
\lqed{1}{6}
\end{proof}



\subsection{The tangent programme and its relation to \texorpdfstring{$E^{(2)}$}{E(2)}}\label{sec:tang2}

\begin{definition}[Tangent programme]\label{def:T}
\begin{equation}\label{eq:T}
 T:=\inf\bigl\{\,g_Q\bigl(\dot\delta+\Aop(\zeta,N_1,Y_1)\bigr)\ :\ (\zeta,N_1,Y_1)
 \ \text{a first-order datum (Def.~\ref{def:data})}\,\bigr\},\qquad
 \kappa_{\rm opt}:=\frac{T}{g_Q(\dot\delta)} .
\end{equation}
By S10 Prop.~4.3, $g_Q(\dot\delta)=\frac12\|u\|_2^2$ and
$g_Q(\delta_c)=s^2g_Q(\dot\delta)(1+o(1))=f_2z^2(1+o(1))$.
\end{definition}

\begin{lemma}[$T$ is a convex programme over a convex cone]\label{lem:convexT}
The set $\mathcal C:=\{\Aop(\zeta,N_1,Y_1)\}\subseteq\mathfrak F$ of admissible directions is
a convex cone, and $\eta\mapsto g_Q(\dot\delta+\eta)$ is a strictly convex quadratic
functional on it; hence $T$ is the squared $g_Q$-distance from $-\dot\delta$ to
$\mathcal C$, every local minimum is global, and the minimiser is unique modulo the kernel
of the form $g_Q$ on $\mathcal C$.  $\mathcal C$ is \emph{not} a linear space
(Thm.~\ref{thm:bdry}, Remark~\ref{rem:coneT0}) and is not claimed closed, so the infimum
need not be attained; this is why $\delta_*$ is defined in \S\ref{sec:kkt} as an orthogonal
residual rather than through a minimiser.
\end{lemma}
\begin{proof}
\lstep{1}{1}{$\mathcal C$ is a convex cone.}
\lby{$\Aop$ is real-linear in $(\zeta,N_1,Y_1)$ (Prop.~\ref{prop:delta1}); the constraint
 set $\{N_1\ge0,\ 0\le Y_1\le N_1\}$ is a convex cone, the set of isometry generators
 $\zeta$ is a convex cone (a convex combination of generators of isometry families
 generates one, by Def.~\ref{def:data}(1) applied to the product family to first order),
 and $\mathfrak F$ is a linear space}
\lstep{1}{2}{Convexity and the distance interpretation.}
\lby{$g_Q$ is a positive quadratic form (Lemma~\ref{lem:norm}), so
 $\eta\mapsto g_Q(\dot\delta+\eta)$ is convex and equals the squared distance from
 $-\dot\delta$ in the Hilbert norm $g_Q^{1/2}$}
\lqed{1}{3}
\end{proof}

\begin{definition}[Hypothesis H1: $g_Q$-differentiability along the scaled family]\label{hyp:H1}
For the first-order datum under consideration,
$g_Q\bigl(\delta(s)-s\,\delta^{(1)}\bigr)=o(s^2)$ as $s\downarrow0$.
\end{definition}

\begin{remark}\label{rem:H1}
H1 is \emph{not} implied by Proposition~\ref{prop:delta1}: the $g_Q$-norm carries the
modular weights $w_{ik}\approx e^{\min(|\kappa_i|,|\kappa_k|)}$ and is strictly stronger
than the Hilbert--Schmidt norm (S8, ``order of limits'').  It is known for the isometric
orbit $(\zeta$ anti-self-adjoint bounded, $N_1=Y_1=0)$: S10 Lemma~6.1$\langle1\rangle3$--%
$\langle1\rangle4$ carries out exactly this estimate, and the non-perturbative check S10
\S7(d) confirms it numerically to $10^{-3}$.  For $N_1,Y_1$ of finite rank with range in a
modular window it is elementary (all remainders are $O(s^2)$ times a fixed element of
$\mathfrak F$).  In general it is a hypothesis.
\end{remark}

\begin{theorem}[Upper bound: the tangent value is achieved by actual channels]\label{thm:upper}
Assume H1 for a first-order datum $(\zeta,N_1,Y_1)$.  Then
\begin{equation}\label{eq:upper}
 \limsup_{s\downarrow0}\ \frac{1}{s^2}\min_{(X,Z)\in\cF}g_Q\bigl(Q-Q_{\rm rec}(X,Z)\bigr)
 \ \le\ g_Q\bigl(\dot\delta+\Aop(\zeta,N_1,Y_1)\bigr).
\end{equation}
If H1 holds for a minimising sequence of data, then
$\limsup_{s\downarrow0}s^{-2}\min_\cF g_Q\le T$, i.e.\
$\min_\cF g_Q\le\kappa_{\rm opt}f_2z^2(1+o(1))$.  Each datum yields an explicit feasible
channel: $X(s),Z(s)$ of \eqref{eq:scaled}, admissible by Lemma~\ref{lem:feas}.
\end{theorem}
\begin{proof}
\lstep{1}{1}{$(X(s),Z(s))\in\cF$ for $0<s<\|N_1\|^{-1}$, so
 $\min_\cF g_Q\le g_Q(\delta(s))$.}
\lby{Lemma~\ref{lem:feas}}
\lstep{1}{2}{$g_Q(\delta(s))^{1/2}\le s\,g_Q(\delta^{(1)})^{1/2}+g_Q(\delta(s)-s\delta^{(1)})^{1/2}
 =s\,g_Q(\delta^{(1)})^{1/2}+o(s)$.}
\lby{the triangle inequality for the Hilbert seminorm $g_Q^{1/2}$ (Lemma~\ref{lem:norm})
 and H1}
\lstep{1}{3}{\eqref{eq:upper}; the statement for a minimising sequence.}
\lby{$\langle1\rangle1$--$\langle1\rangle2$, squaring and dividing by $s^2$; then taking the
 infimum over data, the $\limsup$ being bounded by each individual value}
\lqed{1}{4}
\end{proof}

\begin{corollary}[Consistency with S10]\label{cor:orbit}
Restricting to $N_1=Y_1=0$ and $\zeta=-\mathcal G$ with $\mathcal G$ bounded
anti-self-adjoint and $\mathcal G=E_D\mathcal GE_D$ gives
$\delta^{(1)}=E_I[P,D_\chi+\mathcal G]E_I$, hence
$\inf_{\mathcal G}g_Q(\delta^{(1)})=(1-\theta)g_Q(\dot\delta)$ with $\theta$ of
S10 Thm.~6.2, eq.~(26); therefore $T\le(1-\theta)g_Q(\dot\delta)$, i.e.\
$\kappa_{\rm opt}\le1-\theta$.
\end{corollary}
\begin{proof}
\lstep{1}{1}{$\Aop(-\mathcal G,0,0)=E_I[P,\mathcal G]E_I$.}
\lby{$\zeta=-\mathcal G$, $\zeta^*=\mathcal G$, so \eqref{eq:Ablocks} gives
 $\Aop_{AD}=Q_{AD}\mathcal G$ and
 $\Aop_{DD}=-(\mathcal GQ_{DD}-Q_{DD}\mathcal G)=[Q_{DD},\mathcal G]$; on the other side
 $\mathcal G=E_D\mathcal GE_D$ gives $E_A\mathcal G=\mathcal GE_A=0$, so
 $E_A[P,\mathcal G]E_D=Q_{AD}\mathcal G$, $E_D[P,\mathcal G]E_D=[Q_{DD},\mathcal G]$ and
 $E_A[P,\mathcal G]E_A=0$: the blocks agree}
\lstep{1}{2}{$\delta^{(1)}=E_I[P,D_\chi+\mathcal G]E_I$ and
 $g_Q(\delta^{(1)})=\frac12\|\Pi P^\perp(D_\chi+\mathcal G)P\|_2^2=\frac12\|u+v_{\mathcal G}\|^2$
 with $v_{\mathcal G}=\Pi(P^\perp\mathcal GP)$.}
\lby{$\langle1\rangle1$ and $\dot\delta=E_I[P,D_\chi]E_I$; then S10
 Prop.~4.3$\langle1\rangle2$ (the identity
 $g_Q(E_I[P,\mathcal Z]E_I)=\frac12\|\Pi P^\perp\mathcal ZP\|_2^2$ for anti-self-adjoint
 $\mathcal Z$) and linearity of $\mathcal Z\mapsto\Pi P^\perp\mathcal ZP$}
\lstep{1}{3}{$\inf_{\mathcal G}g_Q(\delta^{(1)})=\frac12\operatorname{dist}(u,\mathcal Y)^2
 =(1-\theta)\frac12\|u\|^2=(1-\theta)g_Q(\dot\delta)$, whence $T\le(1-\theta)g_Q(\dot\delta)$.}
\lby{$\langle1\rangle2$: $\{-v_{\mathcal G}\}$ spans the real subspace $\mathcal Y$ of S10
 Prop.~10.1, $\theta=\|\Pi_{\mathcal Y}u\|^2/\|u\|^2$ by S10 Thm.~6.2; and the infimum in
 \eqref{eq:T} is over a larger set}
\lqed{1}{4}
\end{proof}


\subsection{Isometry directions: boundary-moving generators are singular \texorpdfstring{$N_1$}{N1}-directions}
\label{sec:boundary}

$U_s^*U_s=1$ forces $\zeta+\zeta^*=0$ on $\operatorname{dom}\zeta$, i.e.\ $\zeta$ is
skew-symmetric; but $\zeta$ need not be skew-\emph{adjoint}, and the non-surjective
families (further compression onto a subinterval) are exactly the skew-symmetric,
non-skew-adjoint ones.  They are not a new stratum:

\begin{theorem}[The boundary directions are contained in the $N_1$-cone]\label{thm:bdry}
Let $\zeta$ generate a family of isometries of $\hh_D$ as in Def.~\ref{def:data}(1), let
$\zeta^*$ be its Hilbert-space adjoint.  Define the \emph{boundary form} $\sigma$ and the
anti-symmetric part $\mathcal G$ by
\begin{equation}\label{eq:GS}
 \langle\sigma g,g\rangle:=\Re\langle\zeta^*g,g\rangle\quad(g\in\operatorname{dom}\zeta^*),
 \qquad \mathcal G:=\sigma{-}\zeta^*=-\tfrac12(\zeta^*{-}\zeta^{**})\ \ (\text{anti-symmetric}),
\end{equation}
so that $\zeta^*=-\mathcal G+\sigma$ and $\zeta\subseteq\mathcal G+\sigma$.  (Writing
``$\sigma=\frac12(\zeta+\zeta^*)$'' is \emph{not} well posed: $\zeta$ is not defined on
$\operatorname{dom}\zeta^*$, and on $\operatorname{dom}\zeta$ the symmetric part vanishes
identically --- see (1).  All of $\sigma$'s content sits on
$\operatorname{dom}\zeta^*\setminus\operatorname{dom}\zeta$.)
Then
\begin{enumerate}[leftmargin=1.6em,itemsep=1pt]
\item $\sigma\le0$, and $\sigma=0$ on $\operatorname{dom}\zeta$;
\item as sesquilinear forms, for all $N_1\ge0$, $0\le Y_1\le N_1$,
 \begin{equation}\label{eq:absorb}
  \Aop(\zeta,N_1,Y_1)=\Aop(\mathcal G,\ N_1-2\sigma,\ Y_1),\qquad
  N_1-2\sigma\ \ge\ N_1\ \ge\ Y_1\ \ge\ 0 ;
 \end{equation}
\item hence the \emph{inclusion}
 $\mathcal C\subseteq\{\Aop(\mathcal G,N_1',Y_1):\mathcal G=\sigma{-}\zeta^*\ \text{the
 anti-symmetric part of an admissible}\ \zeta,\ N_1'\ge0\ \text{a form},\ 0\le Y_1\le N_1'\}$
 holds, the extra $N_1'$ being supported on $\partial D$.  (The right-hand side is
 described by \emph{how $\mathcal G$ arises}, not by the property of being itself an
 isometry generator, which $\langle1\rangle4$ does not prove.)  The inclusion is
 \emph{not} an equality and the right-hand side must \emph{not} be enlarged to \emph{all}
 anti-symmetric $\mathcal G$: see Remark~\ref{rem:coneT0};
\item if $\zeta$ is bounded then $\sigma=0$ and $\zeta^*=-\zeta$: bounded isometry families
 give nothing beyond S10's anti-self-adjoint class.
\end{enumerate}
\end{theorem}
\begin{proof}
\lstep{1}{1}{(1).  $U_s^*$ is a contraction semigroup with generator $\zeta^*$, so
 $\Re\langle\zeta^*g,g\rangle\le0$ for $g\in\operatorname{dom}\zeta^*$, i.e.\ $\sigma\le0$;
 and $\|U_sf\|=\|f\|$ gives $\frac{\dd}{\dd s}\|U_sf\|^2|_{0}=2\Re\langle\zeta f,f\rangle=0$
 on $\operatorname{dom}\zeta$.}
\lby{$\|U_s^*\|=\|U_s\|=1$ and the Lumer--Phillips characterisation of contraction
 semigroup generators (dissipativity $\Re\langle\zeta^*g,g\rangle\le0$); $\langle\sigma g,g\rangle=\Re\langle\zeta^*g,g\rangle$
 by \eqref{eq:GS} and $\langle\zeta^{**}g,g\rangle=\overline{\langle \zeta^* g,g\rangle}$ on the common domain}
\lstep{1}{2}{$\zeta=\mathcal G+\sigma$ and $\zeta^*=-\mathcal G+\sigma$, hence
 $\zeta^*Q_{DD}+Q_{DD}\zeta=-[\mathcal G,Q_{DD}]+\{\sigma,Q_{DD}\}$ and
 $Q_{AD}\zeta=Q_{AD}\mathcal G+Q_{AD}\sigma$.}
\lby{\eqref{eq:GS} solved for $\zeta,\zeta^*$; then expand}
\lstep{1}{3}{(2).}
\lby{substituting $\langle1\rangle2$ into \eqref{eq:Ablocks}:
 $\Aop_{DD}=[\mathcal G,Q_{DD}]-\{\sigma,Q_{DD}\}+\frac12\{N_1,Q_{DD}\}-Y_1
 =[\mathcal G,Q_{DD}]+\frac12\{N_1-2\sigma,Q_{DD}\}-Y_1$ and
 $\Aop_{AD}=-Q_{AD}\mathcal G+\frac12Q_{AD}(N_1-2\sigma)$, which are the blocks of
 $\Aop(\mathcal G,N_1-2\sigma,Y_1)$ because $\mathcal G^*=-\mathcal G$; and
 $-2\sigma\ge0$ by $\langle1\rangle1$}
\lstep{1}{4}{(3).}
\lby{$\langle1\rangle3$: every $\Aop(\zeta,N_1,Y_1)$ is an $\Aop(\mathcal G,N_1',Y_1)$ with
 $N_1'=N_1-2\sigma\ge0$ a form and $\mathcal G=\sigma{-}\zeta^*$ anti-symmetric; on the
 right-hand side $\mathcal G$ is read, by definition, as \emph{the anti-symmetric part of an
 admissible} $\zeta$ --- one may not conclude from ``$\mathcal G$ and $\zeta$ agree on
 $C_c^\infty(D)$'' that $\mathcal G$ generates an isometry semigroup, since that needs
 $C_c^\infty(D)$ to be a core for $\zeta$ (true for $\zeta_0,\zeta_\infty$, not assumed in
 general).  No converse is claimed}
\lstep{1}{5}{(4).}
\lby{a bounded symmetric operator is self-adjoint, so
 $\operatorname{dom}\zeta=\operatorname{dom}\zeta^*=\hh_D$ and $\sigma=0$ by
 $\langle1\rangle1$}
\lqed{1}{6}
\end{proof}

\begin{remark}[Example: further compression onto a subinterval]\label{ex:bdry}
Let $\psi\in C^{1,1}(\overline D)$ with $\psi(0)\ge0\ge\psi(1)$ (inward pointing), let
$\varphi_s$ be its flow (which maps $D$ into $D$), and let $U_s:=W_{\varphi_s}$ be the
half-density push-forward $(W_\varphi f)(y)=f(\varphi^{-1}(y))((\varphi^{-1})'(y))^{1/2}$, an isometry of $\hh_D$, onto iff
$\psi(0)=\psi(1)=0$.  Then $\zeta=-D_\psi$, $\zeta^*=D_\psi$ on $H^1(D)$, and
\begin{equation}\label{eq:sigmabd}
 \langle\sigma g,g\rangle=\Re\langle D_\psi g,g\rangle=\tfrac12\bigl[\psi|g|^2\bigr]_0^1
  =\tfrac12\psi(1)|g(1)|^2-\tfrac12\psi(0)|g(0)|^2\ \le\ 0,
\end{equation}
so $-2\sigma=|\psi(1)|\,\delta_1\otimes\delta_1+\psi(0)\,\delta_0\otimes\delta_0\ \ge0$ is a
positive measure concentrated on $\partial D$: the boundary-moving direction is the
$N_1$-direction with $N_1$ singular at the endpoints and $Y_1=0$.  (The forms
$Q_{AD}\sigma$ and $\{\sigma,Q_{DD}\}$ are well defined on $C_c^\infty$ because
$PE_Af_A$ and $PE_Dg$ are real-analytic near $\partial D$ for compactly supported
$f_A,g$.)
\end{remark}

\begin{remark}[Sanity check: over-compression is worse]\label{rem:overcomp}
Take $\psi=\lambda\chi|_D$, $\lambda>0$ (admissible: $\chi(0)=0$, $\chi(1)=-1<0$).  Then
$\Aop_{AD}=\lambda Q_{AD}D_\chi=\lambda\dot\delta_{AD}$ and
$\Aop_{DD}=\lambda[Q_{DD},D_\chi]=0=\lambda\dot\delta_{DD}$ (Lemma~\ref{lem:diagzero}), so
$\delta^{(1)}=(1+\lambda)\dot\delta$ and $g_Q(\delta^{(1)})=(1+\lambda)^2g_Q(\dot\delta)$:
compressing further strictly increases the defect, and the one-sided derivative
$2\lambda g_Q(\dot\delta)>0$ is positive, as the criterion of \S\ref{sec:kkt} requires.
The opposite sign $\lambda<0$ (decompression) is \emph{inadmissible}: it points out of $D$.
\end{remark}


\subsection{The lower bound: what is proved and what is hypothesis}\label{sec:lower}

\begin{proposition}[Coercivity of $g_Q$]\label{prop:coerc}
$w_{ik}\ge1$ for all $i,k$, hence $g_Q(\delta)\ge\frac14\|\delta\|_2^2$.  Consequently any
$(X,Z)\in\cF$ with $g_Q(Q-Q_{\rm rec}(X,Z))\le Cs^2$ satisfies
\begin{equation}\label{eq:HSsmall}
 \|Q_{AD}-Q_{AB}X^*\|_2\le2\sqrt Cs,\qquad \|Q_{DD}-XQ_{BB}X^*-Y\|_2\le2\sqrt Cs .
\end{equation}
\end{proposition}
\begin{proof}
\lstep{1}{1}{$\max\{a(1-b)+b(1-a):a,b\in[0,1]\}=1$.}
\lby{$a(1-b)+b(1-a)=a+b-2ab$ is affine in each variable, so its maximum over the square is
 attained at a vertex, where it takes the values $0,1,1,0$}
\lstep{1}{2}{$w_{ik}\ge1$ and $g_Q(\delta)=\frac14\sum w_{ik}|\widetilde\delta_{ik}|^2\ge\frac14\|\delta\|_2^2$.}
\lby{$\langle1\rangle1$, the definition of $w_{ik}$, and unitary invariance of $\|\cdot\|_2$}
\lstep{1}{3}{\eqref{eq:HSsmall}.}
\lby{$\langle1\rangle2$ applied to the individual blocks of $\delta$ (each block norm is
 bounded by $\|\delta\|_2$) together with \eqref{eq:exactblocks}}
\lqed{1}{4}
\end{proof}

\begin{remark}[Why \eqref{eq:HSsmall} does not localise the minimiser]\label{rem:noloc}
To convert \eqref{eq:HSsmall} into $\|X-X_c\|=O(s)$ one would have to invert $Q_{AB}$ from
the left: $\|Q_{AB}(X_c-X)^*\|_2=O(s)$.  But $Q_{AB}=E_APE_B$ has $0$ in its essential
spectrum (it is a Hankel-type operator between disjoint intervals; its singular values
accumulate at $0$), so no such inversion is available, and the $DD$-constraint only
controls the combination $XQ_{BB}X^*+Y$.  The directions along which $X$ may move by
$O(1)$ without changing $g_Q$ at order $s^2$ are exactly the \emph{invisible} directions
(the kernel of the projection $\Pi$ of S10 Def.~4.2), which Phase~4 showed to be dense in
the exterior directions.  This is the precise obstruction to the lower bound.
\end{remark}

\begin{definition}[Hypothesis H2: localisation of the minimisers]\label{hyp:H2}
There are $C<\infty$, $s_0>0$ such that for $0<s<s_0$ the programme $E^{(2)}=\min_\cF g_Q(Q-Q_{\rm rec})$
has a minimiser $(X_s,Z_s)$ of the form \eqref{eq:scaled} for some first-order datum
$(\zeta^{(s)},N_1^{(s)},Y_1^{(s)})$ with $\|N_1^{(s)}\|+\|Y_1^{(s)}\|\le C$ and
$g_Q(\delta(s)-s\delta^{(1),(s)})=o(s^2)$ uniformly.
\end{definition}

H2 is \emph{not} a regularity assumption: it asserts existence of the minimiser, its scaled
form, boundedness of the data and uniformity of the remainder --- i.e.\ it is the lower
bound in disguise.  It is stated so that everything that depends on it is visible, not
because it is expected to be cheap.

\begin{theorem}[Value of the programme, conditionally]\label{thm:value}
Under H1 (for a minimising sequence of data) and H2,
\begin{equation}\label{eq:valuethm}
 \lim_{s\downarrow0}\ \frac{1}{s^2}\min_{(X,Z)\in\cF}g_Q\bigl(Q-Q_{\rm rec}(X,Z)\bigr)=T,
 \qquad\text{i.e.}\qquad E^{(2)}=\kappa_{\rm opt}\,f_2z^2(1+o(1)),\quad
 \kappa_{\rm opt}=\frac{T}{g_Q(\dot\delta)} .
\end{equation}
\end{theorem}
\begin{proof}
\lstep{1}{1}{$\limsup_s s^{-2}\min_\cF g_Q\le T$.}
\lby{Theorem~\ref{thm:upper}}
\lstep{1}{2}{$\liminf_s s^{-2}\min_\cF g_Q\ge T$.}
\lby{H2 gives minimisers with first-order data; by the triangle inequality for
 $g_Q^{1/2}$ and the uniform remainder,
 $g_Q(\delta(s))^{1/2}\ge s\,g_Q(\delta^{(1),(s)})^{1/2}-o(s)\ge s\,T^{1/2}-o(s)$,
 the last step by the definition \eqref{eq:T} of $T$ as an infimum over all data}
\lqed{1}{3}\lby{$\langle1\rangle1$, $\langle1\rangle2$; $g_Q(\delta_c)=f_2z^2(1+o(1))$ and
 $z=s(1+o(1))$}
\end{proof}

\begin{remark}[An unconditional route to the lower bound]\label{rem:dual}
H2 can be bypassed at the price of an explicit dual witness.  For \emph{every} self-adjoint
trace-class $t$ on $\hh_I$, weak duality (S9 Thm.~4.1, S10~(N1)--(N2)) gives, with
$\Delta_z(t):=\inf_\cF\Tr\,t(Q-Q_{\rm rec})$,
\begin{equation}\label{eq:weakdual}
 \min_{\cF}g_Q\ \ge\ \mathcal D_z(t):=\frac{\Delta_z(t)^2}{8V(t)}\quad(\text{provided }\Delta_z(t)\ge0),\qquad
 \sup_\cF L=2\Re\Tr(t_{DA}Q_{AB}V)+\Tr(t_{DD}V^*R_+V)
  +\Tr\bigl[\bigl((1-V^*V)^{1/2}t_{DD}(1-V^*V)^{1/2}\bigr)_+\bigr],
\end{equation}
the last supremum over $\|V\|\le1$, $V=X^*$.  A lower bound matching \eqref{eq:valuethm}
therefore needs one $t$ (namely $t_*$ of \S\ref{sec:kkt}, at finite $s$) whose multipliers
satisfy the exact sign conditions of S10 Thm.~3.2 at the \emph{optimal} point, not at the
compression.  The criterion of \S\ref{sec:kkt} is precisely the leading order of those
conditions; verifying them at finite $s$ is a task for F2.
\end{remark}

\section{The KKT sign criterion at the isometric optimum}\label{sec:kkt}


Throughout this section $\delta_*$ is the optimal defect \emph{on the linear class}.  We do
\emph{not} posit a minimiser $\mathcal G_*$ (Cor.~\ref{cor:orbit} is an infimum over a class
whose image need not be closed): let
$L:=\{E_I[P,\mathcal G]E_I=[Q,\mathcal G]:\mathcal G\ \text{bounded anti-self-adjoint},\
\mathcal G=E_D\mathcal GE_D\}$ and let $\overline L$ be its closure in the Hilbert space
$(\mathfrak F,b)$; then
\begin{equation}\label{eq:deltastar}
 \delta_*:=\dot\delta-\Pi_{\overline L}\,\dot\delta\ \ (\text{the $b$-orthogonal residual, which always exists}),
 \qquad t_*:=\SLD^{-1}\delta_*,\qquad g_Q(\delta_*)=(1-\theta)g_Q(\dot\delta),
\end{equation}
and the normal equation $b(\delta_*,\eta)=0$ for all $\eta\in\overline L$ holds by
construction.  When a minimiser $\mathcal G_*$ exists,
$\delta_*=\dot\delta+\Aop(-\mathcal G_*,0,0)=E_I[P,D_\chi+\mathcal G_*]E_I$ and
$\zeta_*=-\mathcal G_*$, which is the reading used in \S\ref{sec:f2}.

\begin{definition}[The two criterion operators]\label{def:M}
\begin{equation}\label{eq:Mdef}
 M_*:=E_D\,\Herm(t_*Q)\,E_D=\tfrac12E_D\bigl(t_*Q+Qt_*\bigr)E_D
  =\tfrac12\bigl(t_{*,DD}Q_{DD}+Q_{DD}t_{*,DD}\bigr)
   +\tfrac12\bigl(t_{*,DA}Q_{AD}+Q_{DA}t_{*,AD}\bigr),
 \qquad \tau_*:=t_{*,DD}=E_Dt_*E_D,
\end{equation}
$\Herm(S):=\frac12(S+S^*)$.  \textbf{Caveat (R11).}  $\delta_*\in\mathfrak F$ does
\emph{not} imply $t_*\in\Stwo$, because $\|t_*\|_2^2=\sum w_{ik}^2|\widetilde{(\delta_*)}_{ik}|^2$
while $g_Q(\delta_*)=\frac14\sum w_{ik}|\widetilde{(\delta_*)}_{ik}|^2$ and $w$ is unbounded.
In general $M_*$ and $\tau_*$ are therefore \emph{quadratic forms}, not operators, and
$\Tr(N_1M_*)$ is not a trace; Definition~\ref{def:forms} below gives the intrinsic
definition through $b(\delta_*,\cdot)$, the domain $\mathfrak D$, and the reading of
``$M_*\ge0$'' as a form inequality (then $\lambda_{\min}$ is an infimum over $\mathfrak D$
and may be $-\infty$).  When $t_*$ is trace class the two definitions agree and
$M_*,\tau_*$ are self-adjoint operators on $\hh_D$.
\end{definition}

\begin{lemma}[The one-sided derivative along the non-isometric directions]\label{lem:deriv}
For $N_1\ge0$ and $0\le Y_1\le N_1$ (forms with $\Aop(0,N_1,Y_1)\in\mathfrak F$),
\begin{equation}\label{eq:deriv}
 \frac{\dd}{\dd\sigma}\,g_Q\bigl(\delta_*+\sigma\Aop(0,N_1,Y_1)\bigr)\Big|_{\sigma=0^+}
 =\tfrac12\Tr\bigl(t_*\Aop(0,N_1,Y_1)\bigr)
 =\tfrac12\Bigl[\Tr\bigl(N_1M_*\bigr)-\Tr\bigl(Y_1\tau_*\bigr)\Bigr].
\end{equation}
Along the isometric directions $\Aop(\mathcal G,0,0)$ with $\mathcal G$ in the
\emph{linear} class $\overline L$ of \eqref{eq:deltastar} (bounded anti-self-adjoint,
supported in $D$, and their $b$-limits) the derivative vanishes.  Its vanishing is \emph{not} automatic
along all anti-symmetric $\mathcal G$: for the admissible unbounded generator
$\mathcal G=-\lambda D_\chi|_D$ ($\lambda>0$, Remark~\ref{rem:overcomp}) it equals
$2\lambda(1-\theta)g_Q(\dot\delta)\ge0$, which is $>0$ iff $\theta<1$ --- not proved here (Remark~\ref{rem:whytheta}(i)), numerically $1-\theta=0.616\pm0.003$ (Correction~\ref{corr0924:C2}).  Those directions carry their own one-sided sign
condition, Proposition~\ref{prop:bdryfun} and condition~(e) of
Theorem~\ref{thm:kktmain}.
\end{lemma}
\begin{proof}
\lstep{1}{1}{$\frac{\dd}{\dd\sigma}g_Q(\delta_*+\sigma\eta)|_0=\frac12\Tr(t_*\eta)$.}
\lby{Lemma~\ref{lem:norm}}
\lstep{1}{2}{$\Tr(t_*\Aop)=\Tr(t_{*,AD}\Aop_{DA})+\Tr(t_{*,DA}\Aop_{AD})+\Tr(t_{*,DD}\Aop_{DD})$,
 the $AA$-term vanishing.}
\lby{block multiplication of the trace and $\Aop_{AA}=0$ \eqref{eq:Ablocks}}
\lstep{1}{3}{$\Tr(t_{*,AD}\cdot\frac12N_1Q_{DA})+\Tr(t_{*,DA}\cdot\frac12Q_{AD}N_1)
 =\Tr\bigl(N_1\Herm(t_{*,DA}Q_{AD})\bigr)$.}
\lby{cyclicity gives $\frac12\Tr(N_1Q_{DA}t_{*,AD})+\frac12\Tr(N_1t_{*,DA}Q_{AD})$, and
 $Q_{DA}t_{*,AD}=(t_{*,DA}Q_{AD})^*$ because $t_*=t_*^*$, $Q=Q^*$}
\lstep{1}{4}{$\Tr\bigl(t_{*,DD}(\frac12\{N_1,Q_{DD}\}-Y_1)\bigr)
 =\Tr\bigl(N_1\Herm(t_{*,DD}Q_{DD})\bigr)-\Tr(Y_1\tau_*)$.}
\lby{cyclicity: $\frac12\Tr(t_{*,DD}(N_1Q_{DD}+Q_{DD}N_1))=\frac12\Tr(N_1(Q_{DD}t_{*,DD}+t_{*,DD}Q_{DD}))$}
\lstep{1}{5}{\eqref{eq:deriv}.}
\lby{$\langle1\rangle2$--$\langle1\rangle4$ and $\Herm(t_{*,DA}Q_{AD})+\Herm(t_{*,DD}Q_{DD})
 =\Herm\bigl((t_*Q)_{DD}\bigr)=M_*$, since $(t_*Q)_{DD}=t_{*,DA}Q_{AD}+t_{*,DD}Q_{DD}$}
\lstep{1}{6}{The isometric directions of the linear class.}
\lby{$\mathcal G\mapsto\Aop(\mathcal G,0,0)=[\mathcal G,Q]=-[Q,\mathcal G]$ is real-linear
 with values in $L$ (Corollary~\ref{cor:orbit}$\langle1\rangle1$), and $\delta_*$ is the $b$-orthogonal
 residual of $\dot\delta$ on $\overline L$ \eqref{eq:deltastar}, so
 $b(\delta_*,\eta)=0$ for $\eta\in\overline L$; the counterexample outside $\overline L$
 is Remark~\ref{rem:overcomp} evaluated in Proposition~\ref{prop:bdryfun}(4), which gives $2\lambda(1-\theta)g_Q(\dot\delta)$: a counterexample iff $\theta<1$ (Correction~\ref{corr0924:C2})}
\lqed{1}{7}
\end{proof}

\begin{proposition}[Stationarity on the orbit $=$ condition {\rm(S)} at the new point]\label{prop:S}
$\zeta_*$ is stationary for $g_Q$ along the isometric directions \emph{of the linear class}
$\overline L$ (bounded anti-self-adjoint $\mathcal G=E_D\mathcal GE_D$ and their $b$-limits;
the unbounded ones carry condition~(e) instead) if and only if
\begin{equation}\label{eq:Snew}
 E_D[Q,t_*]E_D=0,\qquad\text{equivalently}\qquad (t_*Q)_{DD}=(t_*Q)_{DD}^*=M_* .
\end{equation}
Thus S10's self-adjointness condition {\rm(S)}, which \emph{fails} at the compression
($E_Du_cE_D\ne0$, S10 Cor.~11.2), holds automatically at the isometric optimum, and there
$M_*$ is not merely the hermitian part of $(t_*Q)_{DD}$ but equals it.
\end{proposition}
\begin{proof}
\lstep{1}{1}{$\Tr\bigl(t_*\Aop(\mathcal G,0,0)\bigr)=\Tr\bigl(E_D[Q,t_*]E_D\,\mathcal G\bigr)$
 for anti-Hermitian $\mathcal G=E_D\mathcal GE_D$.}
\lby{$\Aop(\mathcal G,0,0)=-E_I[P,\mathcal G]E_I=[\mathcal G,Q]$
 (Cor.~\ref{cor:orbit}$\langle1\rangle1$ and $E_I[P,\mathcal G]E_I=[Q,\mathcal G]$ for
 $\mathcal G$ supported in $D$), then cyclicity:
 $\Tr(t_*[\mathcal G,Q])=\Tr(([Q,t_*])\mathcal G)$, and $\mathcal G=E_D\mathcal GE_D$
 projects the first factor onto its $DD$-block}
\lstep{1}{2}{$C:=E_D[Q,t_*]E_D$ is anti-Hermitian, and $\Tr(C\mathcal G)=0$ for all
 anti-Hermitian $\mathcal G$ iff $C=0$.}
\lby{$[Q,t_*]^*=[t_*,Q]=-[Q,t_*]$ and $E_D$ is a projection; for the second claim take
 $\mathcal G=C$, which is an admissible test direction when $C$ is bounded (otherwise
 approximate $C$ by its spectral truncations $\mathbf 1_{[-n,n]}(|C|)C$, which are bounded,
 anti-Hermitian, supported in $D$, and converge to $C$ strongly):
 $\Tr(C^2)=-\Tr(CC^*)=-\|C\|_2^2$}
\lstep{1}{3}{\eqref{eq:Snew} and the identification with $M_*$.}
\lby{$\langle1\rangle1$--$\langle1\rangle2$; $E_DQt_*E_D=((t_*Q)_{DD})^*$ because
 $Q,t_*$ are self-adjoint, so $C=0$ says $(t_*Q)_{DD}$ is self-adjoint, and then
 $\Herm((t_*Q)_{DD})=(t_*Q)_{DD}$ equals $M_*$ by Definition~\ref{def:M}}
\lqed{1}{4}
\end{proof}

\begin{theorem}[KKT sign criterion for the all-channel optimum]\label{thm:kktmain}
Let $\delta_*,t_*,M_*,\tau_*$ be as above (as forms, Def.~\ref{def:forms}) and assume
$\Aop(0,N_1,Y_1)\in\mathfrak F$ for the directions considered.  Then the following are
equivalent:
\begin{enumerate}[leftmargin=2.2em,itemsep=1pt,label=\rm(\alph*)]
\item $M_*\ge0$ \ and \ $M_*-\tau_*\ge0$ (as forms on $\mathfrak D$);
\item[\rm(c)] $\Tr(N_1M_*)\ \ge\ \max_{0\le Y_1\le N_1}\Tr(Y_1\tau_*)
 =\Tr\bigl[\bigl(N_1^{1/2}\tau_*N_1^{1/2}\bigr)_+\bigr]$ for every bounded $N_1\ge0$;
\end{enumerate}
and, \emph{amended after the referee report}, the global statement is
\begin{enumerate}[leftmargin=2.2em,itemsep=1pt,label=\rm(\alph*),start=4]
\item[\rm(d)] $(\zeta_*,0,0)$ is a \emph{global} minimiser of the tangent programme
 \eqref{eq:T} \ $\iff$ \ {\rm(a)} \emph{and} {\rm(e)} hold, where
\item[\rm(e)] \textbf{(endpoint conditions)} $\ \mathcal E(\zeta):=2\,b\bigl(\delta_*,\Aop(\zeta,0,0)\bigr)\ \ge\ 0$
 for every admissible isometry generator $\zeta$ (Def.~\ref{def:data}(1)) with
 $\Aop(\zeta,0,0)\in\mathfrak F$.  \emph{$\mathcal E$ must be evaluated on the blocks
 \eqref{eq:Ablocks}, which contain $\zeta^*$; the commutator form
 $-2b(\delta_*,[Q,\zeta])$ is valid only for bounded anti-self-adjoint $\zeta$ and is blind
 to the boundary form --- see Prop.~\ref{prop:bdryfun}(2$'$).}
\end{enumerate}
Condition~(e) is vacuous on the linear class $\overline L$ (Lemma~\ref{lem:deriv}) and is a
pure \emph{boundary} functional; its explicit form, its reduction to two endpoint numbers,
and the fact that the moving-endpoint one is $\ge0$ (strictly positive iff $\theta<1$, which is not proved; Correction~\ref{corr0924:C2}) are Proposition~\ref{prop:bdryfun}.
Moreover
\begin{enumerate}[leftmargin=2.2em,itemsep=1pt,label=\rm(\alph*),start=2]
\item $M_*\ge(\tau_*)_+$ implies {\rm(a)} and {\rm(c)} (sufficient, strictly stronger than
 {\rm(a)} in general, since the operator order is not a lattice order).
\end{enumerate}
If {\rm(a)} and {\rm(e)} hold then $T=(1-\theta)g_Q(\dot\delta)$, i.e.\
$\kappa_{\rm opt}=1-\theta$ exactly, and (under H1, H2)
$E^{(2)}=(1-\theta)f_2z^2(1+o(1))$; {\rm(a)} alone is necessary but not sufficient.
\end{theorem}


\begin{proof}
\lstep{1}{1}{\lassume{} $0\le Y_1\le N_1$.  \lprove{} $Y_1=N_1^{1/2}CN_1^{1/2}$ with
 $0\le C\le1$, and $\max_{0\le Y_1\le N_1}\Tr(Y_1\tau_*)=\Tr[(N_1^{1/2}\tau_*N_1^{1/2})_+]$.}
\lby{Douglas' lemma: $Y_1\le N_1$ gives $Y_1^{1/2}=N_1^{1/2}K$ with $\|K\|\le1$, so
 $Y_1=N_1^{1/2}KK^*N_1^{1/2}$ and $C:=KK^*\in[0,1]$; then
 $\Tr(Y_1\tau_*)=\Tr(C\,N_1^{1/2}\tau_*N_1^{1/2})$ and
 $\max_{0\le C\le1}\Tr(CS)=\Tr(S_+)$ for $S=S^*$, attained at $C=\mathbf 1_{(0,\infty)}(S)$,
 which corresponds to an admissible $Y_1$}
\lstep{1}{2}{(a) $\Rightarrow$ (c).}
\lby{$\langle1\rangle1$ and, for $0\le Y_1\le N_1$,
 $\Tr(Y_1\tau_*)\le\Tr(Y_1M_*)\le\Tr(N_1M_*)$: the first inequality because
 $M_*-\tau_*\ge0$ and $Y_1\ge0$ (so $\Tr(Y_1(M_*-\tau_*))\ge0$), the second because
 $M_*\ge0$ and $N_1-Y_1\ge0$}
\lstep{1}{3}{(c) $\Rightarrow$ (a).}
\lby{take $N_1=|\psi\rangle\langle\psi|$ with $\|\psi\|=1$: then
 $N_1^{1/2}\tau_*N_1^{1/2}=\langle\psi,\tau_*\psi\rangle|\psi\rangle\langle\psi|$ and (c)
 reads $\langle\psi,M_*\psi\rangle\ge\max\{0,\langle\psi,\tau_*\psi\rangle\}$, i.e.\
 $M_*\ge0$ and $M_*\ge\tau_*$}
\lstep{1}{4}{(d) $\Leftrightarrow$ (a) \emph{and} (e).}
\lby{Lemma~\ref{lem:convexT}: the feasible directions at $\delta_*$ are
 $\overline L+\mathcal C_\zeta+K$ with $\overline L$ the linear class \eqref{eq:deltastar},
 $\mathcal C_\zeta=\{\Aop(\zeta,0,0)\}$ the one-sided cone of admissible isometry
 generators (Thm.~\ref{thm:bdry}(3): an inclusion, \emph{not} an equality --- the cone may
 not be replaced by all anti-symmetric $\mathcal G$, Remark~\ref{rem:coneT0}), and
 $K=\{\Aop(0,N_1,Y_1)\}$; for a differentiable convex function a point is a global
 minimiser iff its derivative is $\ge0$ along every feasible direction; on $\overline L$ it
 vanishes (Lemma~\ref{lem:deriv}$\langle1\rangle6$), on $K$ it is
 $\frac12[\Tr(N_1M_*)-\Tr(Y_1\tau_*)]$, which is $\ge0$ for all admissible $(N_1,Y_1)$ iff
 (c) iff (a) by $\langle1\rangle1$--$\langle1\rangle3$, and on $\mathcal C_\zeta$ it is
 $\mathcal E(\zeta)$ itself (Prop.~\ref{prop:bdryfun}, eq.~\eqref{eq:Efun}), i.e.\ (e)}
\lstep{1}{5}{(b) $\Rightarrow$ (a), and the final claims.}
\lby{$(\tau_*)_+\ge0$ and $(\tau_*)_+\ge\tau_*$ (note (b) says nothing about (e)); if (d) holds then
 $T=g_Q(\delta_*)=(1-\theta)g_Q(\dot\delta)$ by \eqref{eq:deltastar} and
 Corollary~\ref{cor:orbit}, and Theorem~\ref{thm:value} converts this into the statement
 about $E^{(2)}$}
\lqed{1}{6}
\end{proof}

\begin{corollary}[If the criterion fails, the violating direction is a better channel]\label{cor:better}
Suppose $\psi\in\hh_D$, $\|\psi\|=1$, violates {\rm(a)}: either
$\langle\psi,M_*\psi\rangle<0$ (take $N_1:=|\psi\rangle\langle\psi|$, $Y_1:=0$) or
$\langle\psi,(M_*-\tau_*)\psi\rangle<0$ (take $N_1:=Y_1:=|\psi\rangle\langle\psi|$).  Put
$d:=\Aop(0,N_1,Y_1)$ and $\Delta:=\frac14[\Tr(N_1M_*)-\Tr(Y_1\tau_*)]<0$.  Then for
$\sigma_{\rm opt}=-\Delta/g_Q(d)>0$ the first-order datum $(\zeta_*,\sigma_{\rm opt}N_1,
\sigma_{\rm opt}Y_1)$ has
\begin{equation}\label{eq:gain2}
 g_Q\bigl(\delta_*+\sigma_{\rm opt}d\bigr)=g_Q(\delta_*)-\frac{\Delta^2}{g_Q(d)}
 \ <\ (1-\theta)g_Q(\dot\delta),
\end{equation}
and the corresponding channels \eqref{eq:scaled} are feasible for all small $s$
(Lemma~\ref{lem:feas}); they are genuinely non-isometric (for $Y_1=0$: a lossy
$*$-homomorphism into a corner; for $Y_1=N_1$: the maximal noise allowed by the budget).
\textbf{Strictness (R14).}  $\kappa_{\rm opt}<1-\theta$ requires \emph{one} admissible
$\psi\in\mathfrak D$ with $\Delta<0$ \emph{and} $g_Q(d)<\infty$: if the form $M_*$ is
unbounded below only along a sequence $\psi_n$ with $\Delta_n^2/g_Q(d_n)\to0$, then the
infimum is not attained, $T=(1-\theta)g_Q(\dot\delta)$ and $\kappa_{\rm opt}=1-\theta$
after all.  The quantity that decides is $\sup_\psi\Delta^2/g_Q(d)$, as a function of the
grid and of the weight cap (\S\ref{sec:f2}).
\end{corollary}
\begin{proof}
\lstep{1}{1}{$g_Q(\delta_*+\sigma d)=g_Q(\delta_*)+2\sigma\Delta+\sigma^2g_Q(d)$ with
 $\Delta=b(\delta_*,d)=\frac14\Tr(t_*d)$.}
\lby{Lemma~\ref{lem:norm} and \eqref{eq:deriv}}
\lstep{1}{2}{The minimum over $\sigma>0$ is at $\sigma_{\rm opt}=-\Delta/g_Q(d)>0$ with
 value $g_Q(\delta_*)-\Delta^2/g_Q(d)$.}
\lby{$\langle1\rangle1$, a scalar quadratic with $\Delta<0$}
\lstep{1}{3}{Feasibility and the conclusion.}
\lby{Lemma~\ref{lem:feas} ($\sigma_{\rm opt}N_1\ge\sigma_{\rm opt}Y_1\ge0$);
 $g_Q(\delta_*)=(1-\theta)g_Q(\dot\delta)$ and Definition~\ref{def:T}}
\lqed{1}{4}
\end{proof}


\subsection{The same criterion at the compression point, and S10's multipliers}\label{sec:s10cmp}

\begin{proposition}[Identification of the multipliers]\label{prop:s10}
Let $t$ be self-adjoint trace class on $\hh_I$ and let
$\Sigma_1:=t_{DA}Q_{AB}V_c-t_{DD}V_c^*R_-V_c$, $\Sigma_2:=\Sigma_1+t_{DD}
=t_{DA}Q_{AB}V_c+t_{DD}V_c^*R_+V_c$ be S10's multipliers at the compression
($V_c=W$).  Then, as $s\downarrow0$,
\begin{equation}\label{eq:sigmaM}
 \Sigma_2\longrightarrow (tQ)_{DD}=t_{DA}Q_{AD}+t_{DD}Q_{DD},\qquad
 \Sigma_1\longrightarrow (tQ)_{DD}-t_{DD},
\end{equation}
and, \emph{given} the self-adjointness condition {\rm(S)} of S10 Prop.~5.2
($\Sigma_1=\Sigma_1^*$, equivalently $E_DNE_D$ self-adjoint for $N=P^\perp tP$), the
leading-order multipliers are exactly the criterion operators of
Definition~\ref{def:M} evaluated at $\zeta=0$:
\begin{equation}\label{eq:sigmaM2}
 \Sigma_2^{(0)}=\Herm\bigl((tQ)_{DD}\bigr)=M,\qquad \Sigma_1^{(0)}=M-\tau,\qquad \tau=t_{DD}.
\end{equation}
Consequently the three conditions of S10 Thm.~3.2 split as: \;{\rm(S)} $\Sigma_1=\Sigma_1^*$
is the \emph{rotation} ($\zeta$) condition; \;$\Sigma_2\ge0$ is the \emph{loss-of-isometry}
($N_1$) condition $M\ge0$; \;$\Sigma_1\ge0$ is the \emph{noise} ($Y_1$) condition
$M\ge\tau$.
\end{proposition}
\begin{proof}
\lstep{1}{1}{$V_c=W\to1_{\hh_D}$, $R_+=Q_{BB}\to Q_{DD}$, $R_-\to1-Q_{DD}$ and
 $Q_{AB}\to Q_{AD}$ strongly as $s\downarrow0$, whence \eqref{eq:sigmaM}.}
\lby{$B\uparrow D$ as $s\downarrow0$ and $W_s=1-sD_\chi+O(s^2)$ \eqref{eq:Wexp}; the
 second limit follows from $\Sigma_2=\Sigma_1+t_{DD}$}
\lstep{1}{2}{$(tQ)_{DD}=E_Dt PE_D=E_D\tau_{++}E_D+E_DNE_D$ with $\tau_{++}=PtP$,
 $N=P^\perp tP$, which is S10's $\Sigma_2^{(0)}$.}
\lby{$Q=E_IPE_I$ and $E_ItE_I=t$ give $tQ=tPE_I$, hence $E_DtQE_D=E_DtPE_D$; then
 $tP=PtP+P^\perp tP$; this is S10 Prop.~5.2$\langle1\rangle1$}
\lstep{1}{3}{Under (S), $(tQ)_{DD}$ is self-adjoint, so $\Sigma_2^{(0)}=\Herm((tQ)_{DD})=M$
 and $\Sigma_1^{(0)}=M-\tau$.}
\lby{$\langle1\rangle2$ and S10 Prop.~5.2$\langle1\rangle2$: $\Sigma_2^{(0)}$ is
 self-adjoint iff $E_DNE_D$ is, and $E_D\tau_{\pm\pm}E_D$ are self-adjoint automatically;
 for a self-adjoint operator $\Herm$ is the identity}
\lstep{1}{4}{The three-way split.}
\lby{\eqref{eq:sigmaM2} and Theorem~\ref{thm:kktmain}(a); the rotation condition is the
 vanishing of the derivative along the linear space $L$, which at the compression point
 (where $\zeta=0$ is \emph{not} optimal, S10 Thm.~6.2 with Cor.~11.2) fails exactly when
 $E_DuE_D\ne0$}
\lqed{1}{5}
\end{proof}

\begin{remark}[Consistency of the two derivations]\label{rem:consist}
S10 Thm.~3.2$\langle1\rangle7$--$\langle1\rangle8$ derives, from the feasible contraction
family $X_c(1-\frac\eps2\gamma)$ with the freed noise budget used optimally, the exact
condition $\Tr[(\gamma^{1/2}\hat t\gamma^{1/2})_+]\le\Tr(\sigma_2\gamma)$ for all
$\gamma\ge0$ ($\hat t=V_ct_{DD}V_c^*$, $\sigma_2=X_c^*\Sigma_2X_c$).  Conjugating by the
unitary $X_c$ and letting $s\downarrow0$ this is \emph{verbatim} criterion~(c) of
Theorem~\ref{thm:kktmain} with $\gamma\leftrightarrow N_1$, $\hat t\leftrightarrow\tau_*$,
$\sigma_2\leftrightarrow M_*$.  The two derivations are therefore the same computation,
once at the compression ($\zeta=0$, $t=t_c$) and once at the isometric optimum
($\zeta=\zeta_*$, $t=t_*$); the difference is only \emph{which} $t$ enters.  In particular
the S10 certificate already covered the noise directions at the compression point: what was
missing, and is supplied here, is the same statement at the point that is actually optimal
on the isometric orbit.
\end{remark}

\begin{verdictgap}[\;What Theorem~\ref{thm:kktmain} does and does not decide]
It decides the \emph{non-isometric} half of the tangent problem completely: (a) is
necessary and sufficient for the derivative to be $\ge0$ along every
$\Aop(0,N_1,Y_1)$, and it is a pair of form inequalities on two explicitly computable
objects on $\hh_D$.  It does \emph{not} decide the problem by itself: the isometry
directions contribute the one-sided condition (e), which is \emph{not} proved: on the model cone it is
equivalent to $\mathcal E(\zeta_0)\ge0$ (Prop.~\ref{prop:bdryfun}(5); the moving-endpoint term is $\ge0$ by Prop.~\ref{prop:bdryfun}(4)), which holds only numerically, vacuously at $\zeta_0$ (Remark~\ref{rem:T5}); beyond the model cone (e) needs
Hypothesis~\ref{hyp:H3}, under which, read in $\mathfrak F$ and given $\Aop(\zeta_0,0,0)\notin\mathfrak F$ (numerical), (e) holds outright (Remark~\ref{rem:H3space}; Correction~\ref{corr0924:C3}).  Nor does it decide $E^{(2)}$: the passage from the tangent value
$T$ to $\lim s^{-2}\min_\cF g_Q$ uses H1 and H2 (\S\ref{sec:lower}).  The upper bound
$\kappa_{\rm opt}\le1-\theta$ and, if (a) fails, the strict improvement
\eqref{eq:gain2} are unconditional given H1 (Theorem~\ref{thm:upper}).
\end{verdictgap}


\subsection{Amendments after the referee report (2026-09-15)}\label{sec:amend}

\begin{remark}[The isometry-generator cone must not be enlarged]\label{rem:coneT0}
Theorem~\ref{thm:bdry}(3) is an \emph{inclusion}.  Dropping ``the anti-symmetric part of an admissible $\zeta$'' (i.e.\ admitting every anti-symmetric $\mathcal G$)
from the right-hand side is not a harmless relaxation, it destroys the problem.  Take
$\mathcal G:=+\lambda D_\chi|_D$, $\lambda>0$: it is anti-symmetric on $C_c^\infty(D)$ and
supported in $D$, but its flow $x\mapsto x/(1-s\lambda x)$ maps $D=(0,1)$ \emph{onto}
$(0,\frac1{1-s\lambda})\supsetneq D$, so $e^{s\mathcal G}$ is not an operator on $\hh_D$;
consistently, its boundary form is $\Re\langle-\lambda D_\chi g,g\rangle=+\frac\lambda2|g(1)|^2>0$,
violating Theorem~\ref{thm:bdry}(1).  With it one would get $\Aop_{AD}=-\lambda\dot\delta_{AD}$,
$\Aop_{DD}=0$, hence $\delta^{(1)}=(1-\lambda)\dot\delta$ and
$g_Q(\delta^{(1)})=(1-\lambda)^2g_Q(\dot\delta)\to0$ as $\lambda\to1$: the enlarged cone has
tangent value $T=0$.  There is no paradox: for that family $U_s$ is a \emph{contraction},
not an isometry; the corresponding channel has $X^*=$ the restriction $\hh_D\to\hh_B$, whose
defect block is $Q_{AC}$ with $\|Q_{AC}\|_2^2=O(s)$, so $g_Q(\delta(s))=O(s)\ne O(s^2)$ and
H1 fails for it catastrophically.  The one-sided constraint $\sigma\le0$
(Thm.~\ref{thm:bdry}(1)) is exactly what excludes it, and it is the reason the admissible
$\zeta$'s form a cone and not a linear space.  Two further caveats on
Theorem~\ref{thm:bdry}: (i) the absorption identity (2) is an identity of forms that is
\emph{empty} on $C_c^\infty(A)\oplus C_c^\infty(D)$, where $\sigma$ vanishes and
$\mathcal G=\zeta$ --- it is an exact statement about \emph{contraction} families, and for
isometry families it only says where the boundary data can sit; (ii) the singular
$N_1$-directions are \emph{not} in the $g_Q$-closure of the bounded ones: for
$N_1^{(n)}=\lambda n\mathbf 1_{(1-1/n,1)}\rightharpoonup\lambda\delta_1\otimes\delta_1$ one
has $\|Q_{AD}N_1^{(n)}\|_2^2\to\infty$, so $\Aop(0,N_1^{(n)},0)$ diverges in $g_Q$.  They
sit at infinity, and that is why they need their own condition~(e) rather than a limit of
condition~(a).
\end{remark}

\begin{definition}[$M_*$ and $\tau_*$ as quadratic forms]\label{def:forms}
Since $t_*$ need not be Hilbert--Schmidt (Def.~\ref{def:M}, caveat), define directly
\begin{equation}\label{eq:formdef}
 \langle\psi,M_*\psi\rangle:=4\,b\bigl(\delta_*,\Aop(0,|\psi\rangle\langle\psi|,0)\bigr),
 \qquad
 \langle\psi,\tau_*\psi\rangle:=-4\,b\bigl(\delta_*,\Aop(0,0,|\psi\rangle\langle\psi|)\bigr),
\end{equation}
on the domain
$\mathfrak D:=\{\psi\in\hh_D:\Aop(0,|\psi\rangle\langle\psi|,0)\in\mathfrak F\ \text{and}\
\Aop(0,0,|\psi\rangle\langle\psi|)\in\mathfrak F\}$.  Both are finite on $\mathfrak D$ by
Cauchy--Schwarz for $b$, they agree with \eqref{eq:Mdef} whenever $t_*$ is trace class
(Lemma~\ref{lem:norm} and \eqref{eq:deriv}), and ``$M_*\ge0$'', ``$M_*-\tau_*\ge0$'' are
read as form inequalities on $\mathfrak D$.  Accordingly
$\lambda_{\min}(M_*):=\inf\{\langle\psi,M_*\psi\rangle:\psi\in\mathfrak D,\|\psi\|=1\}$,
which may be $-\infty$; the polarisation of \eqref{eq:formdef} extends $M_*,\tau_*$ to
$\Tr(N_1M_*)$, $\Tr(Y_1\tau_*)$ for finite-rank $N_1,Y_1$ with range in $\mathfrak D$, and
by monotone limits to bounded $N_1\ge0$ for which $\Aop(0,N_1,0)\in\mathfrak F$.  Two
conventions, needed for the above to be a \emph{form}: $\mathfrak D$ is to be read as a
linear subspace (equivalently, the polarisation is asserted only for $\psi,\varphi$ with
$\psi\pm\varphi\in\mathfrak D$); $N_1\mapsto\Aop(0,N_1,0)$ is linear, so
$\{N_1:\Aop(0,N_1,0)\in\mathfrak F\}$ is a subspace, but the vector set $\mathfrak D$
defined by the quadratic condition need not be one without this reading.  And the extension
by monotone limits reaches \emph{bounded} $N_1$ only: it does not reach the boundary forms
$-2\sigma_\zeta$ of condition (e) --- e.g.\ $\delta_0\otimes\delta_0$ --- which is why (e)
is a separate condition and not a limiting case of (a) (Remark~\ref{rem:coneT0}(ii)).
\end{definition}


\begin{proposition}[The endpoint functional $\mathcal E$ of condition (e)]\label{prop:bdryfun}
For an admissible isometry generator $\zeta$ (Def.~\ref{def:data}(1)) put, with values in
$[-\infty,+\infty]$,
\begin{equation}\label{eq:Efun}
 \mathcal E(\zeta):=\begin{cases}
 2\,b\bigl(\delta_*,\Aop(\zeta,0,0)\bigr)
  \ \ \bigl(=\tfrac12\Tr(t_*\Aop(\zeta,0,0))\ \text{when $t_*$ is trace class}\bigr),
  & \Aop(\zeta,0,0)\in\mathfrak F,\\[2pt]
 +\infty, & \Aop(\zeta,0,0)\notin\mathfrak F,
 \end{cases}
 \qquad \Aop(\zeta,0,0)\ \text{given by \eqref{eq:Ablocks}} .
\end{equation}
The second case is a convention, not an extension of the form $M_*$ (Def.~\ref{def:forms}
makes $M_*$ finite on $\mathfrak D$): such a $\zeta$ carries \emph{no feasible first-order
datum} (Def.~\ref{def:data}(3) requires $\delta^{(1)}\in\mathfrak F$), so condition~(e) is
\emph{vacuous} there rather than satisfied, and reading $\mathcal E=+\infty$ as a
``margin'' is legitimate only after (a) --- i.e.\ $M_*\ge0$ --- has been granted.
Then:
\begin{enumerate}[leftmargin=1.6em,itemsep=1pt]
\item $\mathcal E$ is additive and positively homogeneous on the cone of admissible
 generators, and $\mathcal E(\zeta)=\mathcal E(\zeta')$ whenever $\zeta-\zeta'\in\overline L$;
 in particular $\mathcal E\equiv0$ on $\overline L$ (Lemma~\ref{lem:deriv}).
\item[{\rm(2$'$)}] \emph{(The commutator form is not admissible here.)}  For bounded
 \emph{anti-self-adjoint} $\zeta$ one has $\Aop(\zeta,0,0)=-[Q,\zeta]$
 (Lemma~\ref{lem:commQ}) and then $\mathcal E(\zeta)=-2b(\delta_*,[Q,\zeta])$.  This is
 \emph{false} for $\zeta$ with a nontrivial boundary form: \eqref{eq:Ablocks} contains
 $\zeta^*$, not $-\zeta$, and replacing $\zeta^*$ by $-\zeta$ deletes exactly the boundary
 term, i.e.\ the whole of $\mathcal E$.  Numerically (F2 T5) the commutator expression
 returns, even for $\zeta_\infty$, the four ladder values $-9.1\cdot10^{-6}$,
 $+1.9\cdot10^{-4}$, $-9.9\cdot10^{-3}$, $+7.4\cdot10^{-2}$ at $h=0.24,0.12,0.06,0.03$, with
 other rungs of the scan reaching $-63$ and $-1.6\cdot10^{6}$: the expression is not merely
 small noise, it is $O(1)$ and worse on the finer scans, and carries no information about
 the true value
 $2\lambda(1-\theta)g_Q(\dot\delta)=1.26$--$1.37\,g_Q(\dot\delta)$.  \textbf{Always evaluate $\mathcal E$ on the
 blocks \eqref{eq:Ablocks}} (or, equivalently, on the boundary form \eqref{eq:Ebdry}: the
 two agree to $2\cdot10^{-10}$--$7.5\cdot10^{-8}$, F2 T5; Correction~\ref{corr0924:C7}).
\item \emph{($\mathcal E$ is a cyclicity anomaly, hence a boundary term.)}  If the
 rearrangement $\Tr(t_*[Q,\zeta])=\Tr([t_*,Q]\zeta)$ is legitimate --- which needs $\zeta$
 bounded (and $t_*$ trace class), so that (2$'$) applies --- then
 $\mathcal E(\zeta)=\frac12\Tr\bigl(E_D[Q,t_*]E_D\,\zeta\bigr)=0$
 by Proposition~\ref{prop:S}.  Thus $\mathcal E$ vanishes identically on the bounded
 generators and is supported on $\partial D$; whenever the absorption identity
 \eqref{eq:absorb} applies with $N_1':=-2\sigma_\zeta$ in the form domain of $M_*$,
 \begin{equation}\label{eq:Ebdry}
  \mathcal E(\zeta)=\tfrac12\Tr\bigl((-2\sigma_\zeta)M_*\bigr)\ \ge0\iff
  \text{$M_*\ge0$ on the boundary form of $\zeta$}.
 \end{equation}
\item \emph{(Two model directions.)}  In the modular frame of \S\ref{sec:modframe}
 ($D=\{y>0\}$, endpoints $y=0$ --- the junction with $A$, i.e.\ $x=0$ --- and $y=+\infty$
 --- the moving endpoint $x=1$) consider the generators of \emph{vector-field} type, $\zeta_v:=-D_v$ with
 $v\in C^{1,1}(\overline D)$ inward.  Modulo $\overline L$ the \emph{vector-field subcone}
 $\{\zeta_v\}$ --- \emph{not}, as far as is proved here, the whole admissible cone; see
 Hypothesis~\ref{hyp:H3} --- is generated by
 \begin{align}
  \zeta_0&:=-E_D\partial_yE_D && (v\equiv1:\ \text{the shift semigroup } (U_ag)(y)=g(y-a),
   \ \text{i.e.\ } x\mapsto1-e^{-a}(1-x)), \label{eq:z0}\\
  \zeta_\infty&:=-\lambda D_\chi|_D && (v=\lambda w,\ w(y)=-4\sinh^2\tfrac y2:\ \text{further
   compression}), \label{eq:zinf}
 \end{align}
 which move the endpoint $y=0$ resp.\ $y=+\infty$ and leave the other fixed.
\item $\mathcal E(\zeta_\infty)=2\lambda(1-\theta)\,g_Q(\dot\delta)\ge0$: the moving-endpoint
 condition $\mathcal E(\zeta_\infty)\ge0$ \textbf{holds}, unconditionally.  \emph{Strict} positivity, $\mathcal E(\zeta_\infty)>0$, is equivalent to $\theta<1$, which is \emph{not} proved (Remark~\ref{rem:whytheta}(i); numerically $1-\theta=0.616\pm0.003$, Remark~\ref{rem:twonum}), and is not claimed here (Correction~\ref{corr0924:C2}).
\item Consequently, on the model cone
 $\{\alpha\zeta_0+\beta\zeta_\infty+\overline L:\alpha,\beta\ge0\}$, condition~(e) is
 equivalent to the single scalar inequality
 \begin{equation}\label{eq:econd}
  \boxed{\ \mathcal E(\zeta_0)=2\,b\bigl(\delta_*,\Aop(\zeta_0,0,0)\bigr)\ \ge\ 0\ },
  \qquad\text{and}\quad \mathcal E(\zeta_0)=\tfrac12\Tr\bigl((-2\sigma_{\zeta_0})M_*\bigr)
  \quad\text{\emph{provided} }-2\sigma_{\zeta_0}\ \text{lies in the form domain of }M_* ,
 \end{equation}
 at the junction endpoint.  (Under Hypothesis~\ref{hyp:H3}, by contrast, and given $\Aop(\zeta_0,0,0)\notin\mathfrak F$ (numerical, see below), (e) holds
 \emph{outright} and \eqref{eq:econd} is not invoked at all --- see
 Remark~\ref{rem:H3space}; \eqref{eq:econd} is then the test for \emph{regularised}
 junction directions.)  The proviso is not automatic and in fact
 \textbf{fails} here: $-2\sigma_{\zeta_0}=\delta_0\otimes\delta_0$ is not a bounded
 operator and $\delta_0\notin\hh_D$, and F2's numerics (Remark~\ref{rem:T5}) indicate
 $\mathcal E(\zeta_0)=+\infty$, i.e.\ the case $\Aop(\zeta_0,0,0)\notin\mathfrak F$ of
 \eqref{eq:Efun}.  So \eqref{eq:econd} is a genuine inequality in $[-\infty,+\infty]$, not
 a finite number; what the numerics establish is its \emph{sign} along every admissible
 regularisation.  This number is not determined by $M_*$'s bulk spectrum and must
 be evaluated separately.
\end{enumerate}
\end{proposition}
\begin{proof}
\lstep{1}{1}{(1).}
\lby{$\Aop(\cdot,0,0)$ is real-linear and $b$ is bilinear; the invariance modulo
 $\overline L$ and the vanishing on $\overline L$ are the normal equation
 $b(\delta_*,\eta)=0$, $\eta\in\overline L$, of \eqref{eq:deltastar}}
\lstep{1}{2}{(2).}
\lby{for bounded anti-self-adjoint $\zeta$, $\Aop(\zeta,0,0)=-[Q,\zeta]$
 (Lemma~\ref{lem:commQ}, whose derivation uses $\zeta^*=-\zeta$ --- this is (2$'$)), and
 $\Tr(t_*[Q,\zeta])=\Tr([t_*,Q]\zeta)=-\Tr(E_D[Q,t_*]E_D\zeta)$ when the rearrangement is
 legitimate, $\zeta=E_D\zeta E_D$; then Proposition~\ref{prop:S}.  For the second claim
 apply \eqref{eq:absorb} and \eqref{eq:deriv} with $N_1=-2\sigma_\zeta$, $Y_1=0$}
\lstep{1}{3}{(3).}
\lby{$U_a$ of \eqref{eq:z0} is an isometry semigroup of $L^2(0,\infty)$ with generator
 $-\partial_y$ on $\{g\in H^1:g(0)=0\}$ and boundary form
 $\langle\sigma g,g\rangle=\frac12[v|g|^2]_0^\infty=-\frac12|g(0)|^2\le0$ concentrated at
 $y=0$; $\zeta_\infty$ has $v=\lambda w$ with $w(0)=w'(0)=0$, so its boundary form is
 concentrated at $y=+\infty$ (Remark~\ref{ex:bdry}); a field vanishing at $y=0$ and
 generating a complete flow inside $D$ has zero boundary form and lies in the class where
 $\mathcal E$ vanishes}
\lstep{1}{4}{(4).}
\lby{Remark~\ref{rem:overcomp}: $\Aop(\zeta_\infty,0,0)=\lambda\dot\delta$; hence
 $\mathcal E(\zeta_\infty)=2\lambda b(\delta_*,\dot\delta)=2\lambda\bigl[g_Q(\delta_*)+b(\delta_*,\Pi_{\overline L}\dot\delta)\bigr]
 =2\lambda g_Q(\delta_*)=2\lambda(1-\theta)g_Q(\dot\delta)$, using
 $\dot\delta=\delta_*+\Pi_{\overline L}\dot\delta$ and $b(\delta_*,\overline L)=0$; hence $\mathcal E(\zeta_\infty)=2\lambda g_Q(\delta_*)\ge0$, because $g_Q$ is a positive form ((C4); equivalently $\theta\le1$); strict positivity would need $\theta<1$, which this step does not give (Correction~\ref{corr0924:C2})}
\lstep{1}{5}{(5).}
\lby{$\langle1\rangle1$, $\langle1\rangle3$, $\langle1\rangle4$ and additivity:
 $\mathcal E(\alpha\zeta_0+\beta\zeta_\infty)=\alpha\mathcal E(\zeta_0)+\beta\mathcal E(\zeta_\infty)$
 with the second term $\ge0$ for $\beta\ge0$ by (4) (only $\ge0$ is used, not strict positivity, Correction~\ref{corr0924:C2}); and $\Aop(\zeta_0,0,0)$ is computed from
 \eqref{eq:Ablocks} with $\zeta_0=-E_D\partial_yE_D$, $\zeta_0^*=E_D\partial_yE_D$ on the
 maximal domain --- \emph{not} from the commutator, by (2$'$)}
\lqed{1}{6}
\end{proof}


\begin{remark}[What F2's T5 endpoint computation says]\label{rem:T5}
In the rigid frame of Proposition~\ref{prop:thetaframe}, with $\mathcal E$ evaluated on the
blocks \eqref{eq:Ablocks} (track F2, \texttt{numerics/optimality\_all/F2\_RESULTS.md} T5,
\texttt{f2\_t5\_endpoint.out}):
\begin{center}\small
\begin{tabular}{@{}lcccc@{}}
\hline
$h$ & $0.24$ & $0.12$ & $0.06$ & $0.03$\\ \hline
$\mathcal E(\zeta_0)/g_Q(\dot\delta)$ (backward, admissible) & $+34.1$ & $+66.9$ & $+131.3$ & $+258.0$\\
$\tfrac12M_*(0,0)/h$ (corner term) & $+10.3$ & $+16.1$ & $+24.7$ & $+37.2$\\
\hline
\end{tabular}
\end{center}
Three things are worth recording.  (i) The \emph{discretisation of $\zeta_0$ must be
one-sided}: the backward (dissipative) difference is the admissible isometry generator and
gives the positive numbers above; the centred difference is anti-self-adjoint, lies in the
linear class, and gives $0$, as Lemma~\ref{lem:deriv} requires; the forward difference is
the \emph{inadmissible} outward direction of Remark~\ref{rem:coneT0} and gives the negative
of the backward value.  The sign of (e) is thus decided by admissibility alone, exactly as
the theory predicts, and the three schemes together are a sharp test of the implementation.
(ii) \emph{What diverges, and which piece is the continuum quantity.}  F2 decomposes the
discrete value as
\begin{equation}\label{eq:T5decomp}
 \mathcal E^{\rm disc}(\zeta_0)=\tfrac12\Bigl[\,\underbrace{M_*(0,0)}_{\text{corner}}
  \;+\;\underbrace{h\,\Tr\bigl((-\Delta)M_*\bigr)}_{\text{upwind viscosity}}
  \;+\;\underbrace{M_*(Y_{\max},Y_{\max})}_{\text{outer cut}}\,\Bigr],
\end{equation}
and the tabulated \emph{total} is dominated by the middle term, which is an artefact of the
one-sided difference: the total grows like $h^{-0.97}$ (ladder ratios $1.960,1.963,1.965$)
whereas the corner term grows like $h^{-0.60}$ (ratios $1.562,1.533,1.508$).  So the total
is \emph{scheme-dependent} and ``driven by the corner term'' would be false; the quantity
that represents the continuum $\mathcal E(\zeta_0)$ is the \textbf{corner row},
$\frac12M_*(0,0)/h=+10.3,+16.1,+24.7,+37.2$, and it is that row, not the total, that should be
quoted for the regularised (finite-$h$) junction directions --- as a divergent value, not as a margin (Correction~\ref{corr0924:C8}).  Both diverge, and the corner row says $M_*(y,y)\to+\infty$ as $y\downarrow0$, i.e.\
$\Aop(\zeta_0,0,0)\notin\mathfrak F$ and $\mathcal E(\zeta_0)=+\infty$ in the sense of the
second case of \eqref{eq:Efun}.  By Definition~\ref{def:data}(3) such a $\zeta$ is not a
feasible first-order datum, so the correct reading is that \textbf{(e) is vacuous at
$\zeta_0$}, not that it holds there with a large margin; what the three schemes do
establish is the \emph{sign} along every admissible regularisation.  Interpreted: moving
the $A$--$D$ junction into $D$ is penalised without bound already at first order, which is
unsurprising --- that direction destroys the $A$--$D$ correlations the $AD$-block of the
defect measures, at a rate no first-order gain can match.  (iii) $\mathcal E(\zeta_\infty)$ reproduces $2\lambda(1-\theta_h)g_Q(\dot\delta)$
to $10^{-15}$, and the $\Aop$-form and the boundary form \eqref{eq:Ebdry} agree to
$2\cdot10^{-10}$--$7.5\cdot10^{-8}$ (Correction~\ref{corr0924:C7}).  \textbf{Conclusion (numerical, not proved):} on the
\emph{model cone} $\{\zeta_0,\zeta_\infty\}$ condition (e) holds --- strictly at $\zeta_\infty$, with a
margin that \emph{decreases} under refinement but stays positive, $\mathcal E(\zeta_\infty)/(\lambda g_Q(\dot\delta))=2(1-\theta_h)=1.37,1.32,1.28,1.26$ at $h=0.24,0.12,0.06,0.03$ (towards $2(1-\theta)=1.232\pm0.006$ by \eqref{eq:thetaF3}; Correction~\ref{corr0924:C5}), and vacuously, not with a margin (the corner value, positive at every $h$, diverges; Correction~\ref{corr0924:C8}),
at $\zeta_0$ --- with no sign change in $h$, in the outer truncation $Y_{\max}$, or in the
frame, and with the three schemes separating exactly as admissibility predicts.  Extending
this to the whole admissible cone is Hypothesis~\ref{hyp:H3}.  Together with (a) (F2 T4)
this gives $\kappa_{\rm opt}=1-\theta=0.616\pm0.003$ conditional on H1, H2 \emph{and} H3.
\end{remark}


\begin{definition}[Hypothesis H3: the cone hypothesis]\label{hyp:H3}
For every admissible isometry generator $\zeta$ (Def.~\ref{def:data}(1)) with
$\Aop(\zeta,0,0)\in\mathfrak F$ there are $\alpha,\beta\ge0$ and $\eta\in\overline L$ with
\begin{equation}\label{eq:H3}
 \Aop(\zeta,0,0)=\alpha\,\Aop(\zeta_0,0,0)+\beta\,\Aop(\zeta_\infty,0,0)+\eta ,
\end{equation}
the identity being read \emph{in $\mathfrak F$} (the convention fixed after
\eqref{eq:formdom}).  This is the whole of H3: there is deliberately no ``equivalent''
formulation in terms of boundary forms alone.  Such a reformulation would be strictly
\emph{weaker} --- for $\zeta=\mathrm im(y)$ of Remark~\ref{rem:H3}(2) one has
$\sigma_\zeta=0$, so a boundary-form version would hold trivially with $\alpha=\beta=0$,
whereas \eqref{eq:H3} then demands $\Aop(\zeta,0,0)\in\overline L$, which is exactly what
may fail --- and the weaker version would not imply (e).  (By the absorption identity
\eqref{eq:absorb} the two readings agree only when
$\Aop(\mathcal G_\zeta,0,0)\in\overline L$, which is unavailable for unbounded
$\mathcal G_\zeta$.)
\end{definition}

\begin{remark}[What H3 asserts once the ambient space is fixed]\label{rem:H3space}
Read in $\mathfrak F$, \eqref{eq:H3} forces $\alpha=0$ \emph{provided} $\Aop(\zeta_0,0,0)\notin\mathfrak F$ (numerical, not proved; Correction~\ref{corr0924:C4}): its left-hand side,
$\Aop(\zeta_\infty,0,0)=\lambda\dot\delta$ and $\eta\in\overline L$ all lie in
$\mathfrak F$, whereas $\Aop(\zeta_0,0,0)\notin\mathfrak F$
(as F2's divergent corner term indicates, Remark~\ref{rem:T5}(ii); Prop.~\ref{prop:bdryfun}(5) proves only that $\delta_0\otimes\delta_0$ is outside the form domain of $M_*$).  So, given this, H3 says precisely:
\emph{every admissible isometry generator that carries a feasible first-order datum is, modulo
$\overline L$, a non-negative multiple of $\zeta_\infty$}; and then, by
Prop.~\ref{prop:bdryfun}(1),(4), $\mathcal E(\zeta)=\beta\,\mathcal E(\zeta_\infty)\ge0$, i.e.\
\textbf{(e) holds outright under H3}, with the junction inequality \eqref{eq:econd} never
invoked.  Two consequences, stated so that nothing is oversold.
\begin{enumerate}[leftmargin=1.6em,itemsep=1pt]
\item The content of H3 is exactly the \emph{exclusion} of the non-vector-field generators
 of Remark~\ref{rem:H3}: it is not a device that makes the junction test operative, and (e)
 should not be advertised as ``collapsing to'' the scalar test \eqref{eq:econd}.
\item \eqref{eq:econd} is nevertheless the right test for \emph{regularised} junction
 directions --- at finite mesh, or for any cut-off $\zeta_0^{(h)}$ with
 $\Aop(\zeta_0^{(h)},0,0)\in\mathfrak F$, the direction \emph{is} feasible and the test is
 non-vacuous.  That is what F2 T5 evaluates (Remark~\ref{rem:T5}), and its positivity is
 the evidence that no admissible regularisation of the junction direction descends.
\end{enumerate}
\end{remark}

\begin{remark}[What H3 excludes, and why it is a hypothesis and not a lemma]\label{rem:H3}
Under H3 read in $\mathfrak F$ (Remark~\ref{rem:H3space}), and given $\Aop(\zeta_0,0,0)\notin\mathfrak F$ (numerical; Correction~\ref{corr0924:C4}), $\alpha=0$ is forced, so
$\mathcal E(\zeta)=\beta\,\mathcal E(\zeta_\infty)\ge0$ by Proposition~\ref{prop:bdryfun}(1)
and condition (e) holds outright; \eqref{eq:econd} remains the operative test only for
regularised junction directions (Remark~\ref{rem:H3space}).  H3 is
\emph{proved} for generators of vector-field type (Prop.~\ref{prop:bdryfun}(3)) and is
\emph{open} in general; two families of admissible generators escape both arguments
available here.
\begin{enumerate}[leftmargin=1.6em,itemsep=1pt]
\item \emph{Shift semigroups of multiplicity $>1$.}  By the Cooper--Wold structure theorem
 every $C_0$-semigroup of isometries is a unitary group $\oplus$ a shift semigroup of some
 multiplicity, and $\hh_D\cong L^2(0,\infty)\otimes\ell^2$ carries shifts of arbitrary
 (including infinite) multiplicity.  Their generators are not first-order differential
 operators and their boundary forms have infinite rank: not a point mass, hence not
 covered by Prop.~\ref{prop:bdryfun}(3).  (The elementary dilation $f\mapsto\sqrt2f(2y)$ is
 \emph{not} the sharp example: it is unitary and is a flow, so it is of vector-field type.)
\item \emph{Unbounded skew-adjoint $\zeta=\mathrm i\,m(y)$}, $m$ real and unbounded,
 generating the unitary group $(U_sg)(y)=e^{\mathrm i sm(y)}g(y)$.  Here $\sigma=0$, so
 \eqref{eq:Ebdry} would give $\mathcal E=0$ --- but \eqref{eq:Ebdry} comes from the
 cyclicity rearrangement, which Prop.~\ref{prop:bdryfun}(2$'$),(2) require $\zeta$ bounded
 for; and the other argument, $\mathcal E\equiv0$ on $\overline L$, applies only if
 $[Q,\mathrm i m]\in\overline L$, i.e.\ only if $\zeta$ is $b$-approximable by bounded
 generators --- exactly what Remark~\ref{rem:coneT0}(ii) warns may fail.  For this family
 $\mathcal E(\zeta)$ is neither computed nor signed here.
\end{enumerate}
Consequently every assertion that ``(e) holds'' in this document means ``(e) holds on the
model cone'' (and, under H3 and given $\Aop(\zeta_0,0,0)\notin\mathfrak F$ (numerical), outright, by Remark~\ref{rem:H3space}), and
$\kappa_{\rm opt}=1-\theta$ is conditional on H1, H2 \emph{and} H3.
Settling H3 --- or evaluating $\mathcal E$ directly on the two families above --- is the
sharpest remaining analytic task of this track.
\end{remark}

\section{What the criterion needs numerically (for F2)}\label{sec:f2}


Everything in Theorem~\ref{thm:kktmain} is explicit once $t_*$ is known.  We list the exact
objects F2 has to evaluate, in the $E_I$-visible language and in the form in which the
spectral circle model (E2's \texttt{e2\_common.py}) and the modular Galerkin model
(\texttt{kkt\_continuum.py}) produce them.

\begin{enumerate}[leftmargin=1.8em,itemsep=2pt]
\item \textbf{The optimal isometric defect.}  With $\mathcal G$ anti-Hermitian supported in
 $D$ one has $E_I[P,\mathcal G]E_I=[Q,\mathcal G]$ (S9 \S5.2), so
 \begin{equation}\label{eq:f2delta}
  \delta_*=\dot\delta+[Q,\mathcal G_*],\qquad
  \mathcal G_*=\arg\min_{\mathcal G^*=-\mathcal G,\ \mathcal G=E_D\mathcal GE_D}
   g_Q\bigl(\dot\delta+[Q,\mathcal G]\bigr),
 \end{equation}
 i.e.\ $\mathcal G_*$ is the conjugate-gradient solution already computed for $\theta$
 (S10 \S7(c), S9 \S5.2: $\mathfrak A\mathcal G_*=\beta$).  Normalisation check:
 $g_Q(\delta_*)=(1-\theta)g_Q(\dot\delta)$.
\item \textbf{The SLD.}  In an eigenbasis $\{e_i\}$ of $Q=Q_{II}$ ($Qe_i=q_ie_i$,
 $q_i=(1+e^{\kappa_i})^{-1}$),
 \begin{equation}\label{eq:f2sld}
  \widetilde{(t_*)}_{ik}=w_{ik}\,\widetilde{(\delta_*)}_{ik},\qquad
  w_{ik}=\frac1{q_i(1-q_k)+q_k(1-q_i)}=1+\frac{\cosh\frac{\kappa_i+\kappa_k}2}{\cosh\frac{\kappa_i-\kappa_k}2},
 \end{equation}
 i.e.\ $t_*=\SLD^{-1}\delta_*$ is the entrywise product of the weight matrix with the
 defect in the $Q$-eigenbasis.  Check: $\frac14\Tr(t_*\delta_*)=g_Q(\delta_*)$ and
 $V(t_*)=\Tr(t_*Qt_*(1-Q))=2g_Q(\delta_*)$ (Lemma~\ref{lem:norm}).
\item \textbf{The two criterion operators} (both on $\hh_D$, both self-adjoint):
 \begin{equation}\label{eq:f2M}
  \boxed{\ \tau_*=E_D\,t_*\,E_D,\qquad
  M_*=\tfrac12E_D\bigl(t_*Q+Qt_*\bigr)E_D=\tfrac12E_D\bigl(t_*P+Pt_*\bigr)E_D\ }
 \end{equation}
 (the second form uses $Q=E_IPE_I$ and $E_It_*=t_*$; in a discretisation where $Q$ is the
 primary object the first form is the one to use).  In blocks,
 $M_*=\Herm(t_{*,DA}Q_{AD})+\Herm(t_{*,DD}Q_{DD})$: the $DA$ block of $t_*$ is needed, not
 only the $DD$ block.
\item \textbf{The test, part 1: condition (a).}  Report $\lambda_{\min}(M_*)$ and
 $\lambda_{\min}(M_*-\tau_*)$ in units of $g_Q(\dot\delta)$ (and state which SLD
 normalisation is used, Remark~\ref{rem:factor2}).  No search over general $N_1\ge0$ is
 needed: criterion~(c) is \emph{equivalent} to the two rank-one inequalities.  Because
 $M_*,\tau_*$ are forms (Def.~\ref{def:forms}), a negative $\lambda_{\min}$ is
 \emph{not} by itself a better channel: report also
 \begin{equation}\label{eq:gainreport}
  \sup_{\psi}\ \frac{\Delta(\psi)^2}{g_Q(d_\psi)}\Big/g_Q(\dot\delta),\qquad
  \Delta(\psi)=\tfrac14\bigl[\langle\psi,M_*\psi\rangle-\langle\psi,\tau_*\psi\rangle_{\rm(if\ }Y_1=N_1)\bigr],
  \quad d_\psi=\Aop(0,|\psi\rangle\langle\psi|,Y_1),
 \end{equation}
 as a function of the grid \emph{and} of the weight cap $\kappa_c$: only a limit bounded
 away from $0$ establishes $\kappa_{\rm opt}<1-\theta$ (Cor.~\ref{cor:better}, R14).  A
 violation that scales to zero with the cap is a modular-ultraviolet artefact of the
 regularisation, not a channel.
\item \textbf{The test, part 2: condition (e) --- the endpoint test.}  Evaluate the single
 number \eqref{eq:econd},
 \[ \mathcal E(\zeta_0)/g_Q(\dot\delta),\qquad
   \mathcal E(\zeta_0)=2\,b\bigl(\delta_*,\Aop(\zeta_0,0,0)\bigr)
   =\tfrac12\Tr\bigl(t_*\Aop(\zeta_0,0,0)\bigr),\qquad \zeta_0=-E_D\,\partial_y\,E_D , \]
 in the rigid parameter-free frame of Proposition~\ref{prop:thetaframe} ($Q$ the multiplier,
 $E_D=\mathbf 1_{y>0}$: a matrix computation, no geometry).  \textbf{Use the blocks
 \eqref{eq:Ablocks}} --- $\Aop_{AD}=-Q_{AD}\zeta_0$,
 $\Aop_{DD}=-(\zeta_0^*Q_{DD}+Q_{DD}\zeta_0)$, with $\zeta_0^*$ the transpose of the
 one-sided difference --- \emph{not} the commutator $-[Q,\zeta_0]$, which is blind to the
 boundary form and carries no information about $\mathcal E$: it is not merely solver noise but $O(1)$ and worse on the finer scans (Prop.~\ref{prop:bdryfun}(2$'$); Correction~\ref{corr0924:C6}); and
 discretise $\partial_y$ by the \emph{backward} (dissipative) difference, the centred one
 being the anti-self-adjoint member of the linear class and the forward one inadmissible
 (Remark~\ref{rem:T5}).  Two cross-checks: (i) the analytic value along the other model
 direction,
 $\mathcal E(\zeta_\infty)=2\lambda(1-\theta)g_Q(\dot\delta)$ with
 $1-\theta=0.616\pm0.003$ (Prop.~\ref{prop:bdryfun}(4)) --- the numerics must reproduce it;
 (ii) the ``cell'' form of \eqref{eq:Ebdry}: the diagonal entry of $M_*$ on the first cell
 at $y=0$ (resp.\ the last cell, $y=Y_{\max}$), suitably normalised by the cell width, must
 have the same sign as $\mathcal E(\zeta_0)$ (resp.\ $\mathcal E(\zeta_\infty)$).  Report
 the $h$- and $Y_{\max}$-dependence, and report the \emph{corner term} separately from the
 total: by \eqref{eq:T5decomp} the total carries a scheme-dependent $O(h)$ viscosity piece,
 and only the corner term represents the continuum value.  This tests (e) on the model cone
 only; the non-vector-field generators of Remark~\ref{rem:H3} are not reached by any of
 these directions, and excluding them is exactly the content of H3
 (Remark~\ref{rem:H3space}).
\item \textbf{Stability requirements.}  (i) $L$-stability (grid size) and stability in the
 weight cap $w\mapsto\min(w,1+\cosh\kappa_c)$ of S8 (the cap is a lower bound for $g_Q$ and
 changes $t_*$, so the sign test must be reported as a function of $\kappa_c$);
 (ii) stability in the ultraviolet taper of the spectral model (S10 \S10 Remark: the taper
 removes the region where $\mathcal G_*$ lives); (iii) the same test in the modular
 Galerkin frame, where no cap and no taper are used, as the arbiter.
\item \textbf{A by-product worth reporting.}  $\Herm((t_*Q)_{DD})=M_*$ and
 $\mathrm{AntiHerm}((t_*Q)_{DD})$: the latter is the leading-order violation of S10's
 condition {\rm(S)} \emph{at the new point}, and it must vanish if $\mathcal G_*$ is
 optimal on the orbit (Lemma~\ref{lem:deriv}$\langle1\rangle6$).  This is a free
 consistency check on the numerics: $\|\mathrm{AntiHerm}((t_*Q)_{DD})\|$ should be
 $0$ to solver tolerance, whereas at the compression it equals
 $\|E_Du_cE_D\|\ne0$ (S10 Cor.~11.2).  It is also the \emph{M\"obius-rigidity arbiter}:
 the circle model violates $E_D\dot\delta E_D=0$ (Lemma~\ref{lem:diagzero}) at $O(1/L)$, and
 Prop.~\ref{prop:delta1}$\langle1\rangle5$ uses that identity, so this norm must be reported
 alongside the two eigenvalues and $\mathcal E(\zeta_0)$.
\end{enumerate}

\section{A closed form for \texorpdfstring{$\theta$}{theta}: setup and status}\label{sec:theta}


\subsection{A frame-free form of the whole first-order problem}\label{sec:framefree}

\begin{lemma}[The compression defect is a commutator with $Q$]\label{lem:commQ}
Let $d:=E_ID_\chi E_I$.  Then $\dot\delta=[Q,d]$, and consequently the entire tangent
problem is expressed in terms of the pair $(Q,E_D)$ and the single operator $d$:
\begin{equation}\label{eq:framefree}
 \delta^{(1)}(\zeta,0,0)=[Q,\,d-\zeta]\qquad(\zeta^*=-\zeta,\ \zeta=E_D\zeta E_D),
\end{equation}
the $N_1$- and $Y_1$-terms being those of \eqref{eq:Ablocks}.
\end{lemma}
\begin{proof}
\lstep{1}{1}{$(1-E_I)D_\chi E_I=0$ and $E_ID_\chi(1-E_I)=0$.}
\lby{$D_\chi$ is local and $\chi$ is $C^{1,1}$ with $\chi(0)=\chi'(0)=0$, so $D_\chi f$ is
 supported in $\operatorname{supp}f$; the second identity is the adjoint of the first
 ($D_\chi^*=-D_\chi$)}
\lstep{1}{2}{$E_I[P,D_\chi]E_I=E_IPE_ID_\chi E_I-E_ID_\chi E_IPE_I=Qd-dQ$.}
\lby{insert $1=E_I+(1-E_I)$ between $P$ and $D_\chi$ in both terms and use
 $\langle1\rangle1$ to kill the two cross terms}
\lstep{1}{3}{The second claim.}
\lby{$\langle1\rangle2$; and $\Aop(\zeta,0,0)=-E_I[P,\zeta]E_I=-[Q,\zeta]$ for
 anti-symmetric $\zeta=E_D\zeta E_D$ by Cor.~\ref{cor:orbit}$\langle1\rangle1$ with
 $\mathcal G=-\zeta$, using $E_I[P,\mathcal G]E_I=[Q,\mathcal G]$ for $\mathcal G$
 supported in $D$ (S9 \S5.2)}
\lqed{1}{4}
\end{proof}

\subsection{The modular frame of $I$}\label{sec:modframe}

$I=(-\infty,1)$ is a \emph{half-line}, so by Bisognano--Wichmann (chiral version: the
vacuum modular group of a half-line is the dilation group fixing its endpoint; Hislop--Longo)
the modular coordinate
\begin{equation}\label{eq:ycoord}
 y:=-\log(1-x)\ :\ I\xrightarrow{\ \cong\ }\mathbb R,\qquad
 A=(-\infty,0)\mapsto\{y<0\},\quad D=(0,1)\mapsto\{y>0\},
\end{equation}
implemented unitarily on $L^2$ by the half-density push-forward, turns the vacuum symbol
into a \emph{Fourier multiplier}
\begin{equation}\label{eq:Qmult}
 Q=\bigl(1+e^{-2\pi\mathsf p}\bigr)^{-1},\qquad \mathsf p=-i\partial_y,\qquad
 \text{i.e.}\quad \kappa=-2\pi p\ \ \text{(sign a convention)},
\end{equation}
while $E_D$ becomes the sharp cutoff $\mathbf 1_{\{y>0\}}$ and the compression field becomes
\begin{equation}\label{eq:wfield}
 d=D_w,\qquad w(y)=\frac{\chi(x)}{1-x}\Big|_{x=1-e^{-y}}=-4\sinh^2\!\tfrac y2\ \mathbf 1_{\{y>0\}} ,
\end{equation}
which is $C^{1,1}$ at $y=0$ ($w\sim-y^2$), vanishes identically on $A$, and grows like
$-e^{y}$ as $y\to+\infty$ --- the moving endpoint $x=1$ has been sent to $y=+\infty$.

\begin{proposition}[The $\theta$-problem in the modular frame]\label{prop:thetaframe}
In the frame \eqref{eq:ycoord}--\eqref{eq:wfield},
\begin{equation}\label{eq:thetaWH}
 1-\theta=\frac{\inf_{\mathcal G}g_Q\bigl([Q,d+\mathcal G]\bigr)}{g_Q\bigl([Q,d]\bigr)},
 \qquad \mathcal G\ \text{bounded anti-self-adjoint},\ \
 \mathcal G=\mathbf 1_{\{y>0\}}\mathcal G\mathbf 1_{\{y>0\}}\ \ (\text{the class of Cor.~\ref{cor:orbit}}),
\end{equation}
and, writing $\widehat X(p,p')$ for the kernel of $X$ in the Fourier variable of $y$,
\begin{equation}\label{eq:kernelK}
 g_Q\bigl([Q,X]\bigr)=\tfrac14\iint K(p,p')\,\bigl|\widehat X(p,p')\bigr|^2\dd p\dd p',
 \qquad
 K(p,p')=\frac{\sinh^2\bigl(\pi(p-p')\bigr)}
 {\cosh\bigl(\pi(p-p')\bigr)\,\bigl[\cosh\bigl(\pi(p+p')\bigr)+\cosh\bigl(\pi(p-p')\bigr)\bigr]} .
\end{equation}
Thus $\theta$ is the value of an explicit quadratic Wiener--Hopf problem: minimise a
translation-\emph{non}-invariant quadratic form with kernel $K$ over kernels supported in
the quadrant $\{y>0,y'>0\}$.
\end{proposition}
\begin{proof}
\lstep{1}{1}{\eqref{eq:thetaWH}.}
\lby{Lemma~\ref{lem:commQ} and Corollary~\ref{cor:orbit} (the isometric orbit contributes
 $[Q,\mathcal G]$), divided by $g_Q(\dot\delta)$; unitary equivalence of frames leaves
 $g_Q$ invariant because $g_Q$ is defined by the spectral data of $Q$ alone}
\lstep{1}{2}{With $q_i=(1+e^{\kappa_i})^{-1}$,
 $\widetilde{[Q,X]}_{ik}=(q_i-q_k)\widetilde X_{ik}$ and
 $w_{ik}(q_i-q_k)^2=\dfrac{\sinh^2\frac{\kappa_i-\kappa_k}2}
 {\cosh\frac{\kappa_i-\kappa_k}2\bigl[\cosh\frac{\kappa_i+\kappa_k}2+\cosh\frac{\kappa_i-\kappa_k}2\bigr]}$.}
\lby{with $a=\kappa_i$, $b=\kappa_k$, $s=\frac{a+b}2$, $t=\frac{a-b}2$:
 $q_i-q_k=\frac{e^b-e^a}{(1+e^a)(1+e^b)}=\frac{-2e^s\sinh t}{2e^s(\cosh s+\cosh t)}$ and
 $q_i(1-q_k)+q_k(1-q_i)=\frac{e^a+e^b}{(1+e^a)(1+e^b)}=\frac{\cosh t}{\cosh s+\cosh t}$;
 divide}
\lstep{1}{3}{\eqref{eq:kernelK}.}
\lby{$\langle1\rangle2$ with $\kappa=-2\pi p$ \eqref{eq:Qmult}, so
 $\frac{\kappa_i-\kappa_k}2=-\pi(p-p')$ and $\frac{\kappa_i+\kappa_k}2=-\pi(p+p')$, both
 entering through even functions; and \eqref{eq:gQ}}
\lqed{1}{4}
\end{proof}


\subsection{What $\theta$ measures, and the status of a closed form}\label{sec:thetastatus}

\begin{remark}[The obstruction lives at the moving endpoint]\label{rem:whytheta}
In the frame \eqref{eq:ycoord} the competitor class is $\{\mathcal G\}$: anti-Hermitian and
supported in $\{y>0\}$; the target is $d=D_w$, which \emph{is} supported in $\{y\ge0\}$ but
is only skew-\emph{symmetric}: by \eqref{eq:sigmabd} its symmetric part is the boundary form
at $y=+\infty$ (equivalently at $x=1$), the endpoint the compression moves.  If $d$ were
anti-self-adjoint, $\mathcal G=-d$ would be admissible and $\theta$ would be $1$.  Hence
\begin{equation}\label{eq:thetameaning}
 1-\theta=\ \text{the $g_Q$-distance}^2\ \text{from $[Q,d]$ to $\{[Q,\mathcal G]\}$}\,/\,g_Q([Q,d])
 \ =\ \text{the irreducible ``endpoint'' part of the compression defect} .
\end{equation}
Two consequences.  (i) $\theta<1$ iff $[Q,d]\notin\overline{\{[Q,\mathcal G]\}}$, which is
exactly the statement that the endpoint part cannot be undone by an interior rotation; no
proof of $\theta<1$ from first principles is recorded, although $\theta\le1$ is trivial and
$1-\theta=0.616\pm0.003$ numerically after reconciliation (Remark~\ref{rem:twonum}).  (ii) $\theta>0$ iff $E_DuE_D\ne0$ (S10
Thm.~6.2 with Cor.~11.2): $E_DuE_D$ is purely anti-self-adjoint by M\"obius rigidity, so any
nonzero $DD$-part of $u$ is automatically an admissible descent direction.
\end{remark}

\begin{remark}[The reconciled value of $\theta$ (track F3), superseding the earlier window]\label{rem:twonum}
This remark previously argued that the spectral and Galerkin values of $\theta$ could not
agree before the mesh resolved the region $y\gg1$ near the moving endpoint.  \textbf{That
diagnosis is refuted and the window is superseded.}  Track F3
(\texttt{numerics/optimality\_all/f3\_t0.py}, \texttt{F3\_RESULTS.md}) computes $\theta$ in
the modular Galerkin frame \emph{and} by a direct discretisation of the rigid frame of
Proposition~\ref{prop:thetaframe}; the two agree to $3\times10^{-4}$ after extrapolation
and give
\begin{equation}\label{eq:thetaF3}
 \theta=0.384\pm0.003,\qquad 1-\theta=0.616\pm0.003
\end{equation}
(measured convergence order $h^{0.77}$, five mesh families; the modular window $\Lambda$ is
irrelevant).  The residual $h$-dependence is a \emph{bulk} effect --- the endpoint
contribution saturates already at coarse $h$ --- i.e.\ the opposite of the mechanism
conjectured here.  S10's $0.300\pm0.005$ is not a competing estimate but a
\emph{taper-suppressed lower bound}: relaxing the circle model's smoothstep ultraviolet
taper moves $\theta$ at $L=256$ from $0.2896$ to $0.36$ and beyond.  The one statement of
the old remark that survives is the behaviour of the kernel: $K(p,p')\to1$ as
$|p-p'|\to\infty$ at fixed $p+p'$, while $K\to0$ as $|p+p'|\to\infty$ at fixed $p-p'$ ---
and the latter is precisely the direction the circle taper suppresses, consistent with the
downward bias.  Consequently $\kappa_{\rm opt}=1-\theta=0.616\pm0.003$ if the criterion of
Theorem~\ref{thm:kktmain} holds.
\end{remark}

\begin{verdictgap}[\;$\theta$ is now numerically settled; a closed form is still open]
Proposition~\ref{prop:thetaframe} reduces $\theta$ to an explicit, parameter-free
Wiener--Hopf variational problem: all geometry has been scaled out (no $a,b,c,L,z$ appear
in \eqref{eq:thetaWH}--\eqref{eq:kernelK}), so $\theta$ \emph{is} a pure number --- which
the reconciled numerics \eqref{eq:thetaF3} confirm.  \textbf{We have not solved it}, and no
constant is selected.  The three candidates circulated in the Phase-5 brief are now
\emph{excluded} by \eqref{eq:thetaF3}: $3/\pi^2=0.30396$ and $1/3$ lie far below, and
$1-2/\pi=0.36338$ --- the closest --- is still $6.9\sigma$ low.  That all three fell inside
the earlier $0.30$--$0.40$ window is a cautionary datum about unmotivated candidates, not
evidence for any of them.  What is proved about $\theta$: it is the value of
\eqref{eq:thetaWH}; it is geometry-independent; $\theta>0$ iff $E_DuE_D\ne0$ (S10 Thm.~6.2
with Cor.~11.2), the non-vanishing being numerical.  What is \emph{not} proved: $\theta<1$
(only the converse implication of Remark~\ref{rem:whytheta}(i) is established, namely that
$\theta=1$ would force $[Q,d]\in\overline{\{[Q,\mathcal G]\}}$); the numerics give
$1-\theta=0.616\pm0.003$.
\end{verdictgap}

\section{Summary: proved, conditional, and to be supplied}\label{sec:summary}


\paragraph{Proved here, unconditionally.}
\begin{enumerate}[leftmargin=1.8em,itemsep=1pt]
\item \emph{Exact chart of $\cF$} (Lemmas~\ref{lem:NY}--\ref{lem:WU},
 Prop.~\ref{prop:blocks}): $(X,Z)\in\cF\iff N=1-XX^*\ge0$ and $0\le Y=Z-XQ_{BB}X^*\le N$;
 $X=(1-N)^{1/2}V$ with the \emph{polar} $V$ a co-isometry iff $1\notin\sigma_p(N)$ (sharp,
 Remark~\ref{rem:shift}; \emph{some} co-isometric factor exists iff
 $\dim\ker(1-N)\le\dim\ker X$, Lemma~\ref{lem:polar}(4)), and $V^*=WU$ with $U$ an isometry
 of $\hh_D$.  All defect blocks exactly, no expansion.
\item \emph{First-order defect} (Prop.~\ref{prop:delta1}), with all signs:
 $\delta^{(1)}=\dot\delta+\Aop(\zeta,N_1,Y_1)$, $\Aop_{AD}=-Q_{AD}\zeta+\frac12Q_{AD}N_1$,
 $\Aop_{DD}=-(\zeta^*Q_{DD}+Q_{DD}\zeta)+\frac12\{N_1,Q_{DD}\}-Y_1$; $\dot\delta$ is purely
 $A$--$D$ off-diagonal (Lemma~\ref{lem:diagzero}) and $\dot\delta=[Q,d]$
 (Lemma~\ref{lem:commQ}).
\item \emph{Feasibility of the scaled family} (Lemma~\ref{lem:feas}): each first-order
 datum gives \emph{exactly} feasible channels at every small $s$ --- so the tangent value
 is an upper bound (Thm.~\ref{thm:upper}, given H1), and any violating direction is a
 better channel (Cor.~\ref{cor:better}).
\item \emph{Boundary-moving isometries} (Thm.~\ref{thm:bdry}, as amended): \emph{bounded}
 isometry generators give nothing beyond S10's anti-self-adjoint class (part (4), the part
 that survives intact).  For unbounded ones the absorption identity
 $\Aop(\zeta,N_1,Y_1)=\Aop(\mathcal G,N_1-2\sigma,Y_1)$ with $-2\sigma\ge0$ supported on
 $\partial D$ gives an \emph{inclusion} of cones, not an equality: on the right-hand side
 $\mathcal G$ is the anti-symmetric \emph{part of an admissible} $\zeta$ (that it itself
 generates an isometry semigroup is not claimed), and the singular $N_1$'s are not in the
 $g_Q$-closure of the bounded ones.  Enlarging the cone to all anti-symmetric
 $\mathcal G$ would give $T=0$ (Remark~\ref{rem:coneT0}).  So the answer to the Phase-5
 brief's question is: the boundary directions do not enlarge the \emph{value} through
 $N_1$, but they do form a genuine extra one-sided cone, which is why the criterion needs
 condition (e).
\item \emph{KKT sign criterion} (Thm.~\ref{thm:kktmain}) with the operators
 $M_*=\frac12E_D(t_*Q+Qt_*)E_D$ and $\tau_*=E_Dt_*E_D$:
 the derivative condition along \emph{the non-isometric directions} $\Aop(0,N_1,Y_1)$ is
 (a) $M_*\ge0$ and $M_*\ge\tau_*$, equivalently (c)
 $\Tr(N_1M_*)\ge\Tr[(N_1^{1/2}\tau_*N_1^{1/2})_+]$ for all $N_1\ge0$ --- with $M_*,\tau_*$
 read as quadratic forms (Def.~\ref{def:forms}); $M_*\ge(\tau_*)_+$ is sufficient but
 strictly stronger.  \textbf{Global} optimality of the tangent problem is (a)
 \emph{together with} the endpoint condition (e)
 $\mathcal E(\zeta):=2b(\delta_*,\Aop(\zeta,0,0))\ge0$ along the admissible isometry
 generators --- evaluated on the blocks \eqref{eq:Ablocks}, \emph{not} on $-[Q,\zeta]$,
 which is blind to the boundary form (Prop.~\ref{prop:bdryfun}(2$'$)) --- a pure boundary
 functional which reduces, \emph{on the model cone}, to the single inequality
 $\mathcal E(\zeta_0)\ge0$ at the junction endpoint; the moving-endpoint half is proved:
 $\mathcal E(\zeta_\infty)=2\lambda(1-\theta)g_Q(\dot\delta)\ge0$, strictly positive iff $\theta<1$, which is \emph{not} proved (Remark~\ref{rem:whytheta}(i); numerically $1-\theta=0.616\pm0.003$; Correction~\ref{corr0924:C2}).  Under
 Hypothesis~\ref{hyp:H3}, read in $\mathfrak F$, (e) holds \emph{outright} if
 $\Aop(\zeta_0,0,0)\notin\mathfrak F$ (which is numerical, not proved: F2's divergent corner term, Remark~\ref{rem:T5}(ii); Correction~\ref{corr0924:C4}), for then \eqref{eq:H3} forces $\alpha=0$ and only
 $\mathcal E(\zeta_\infty)\ge0$ is left (Remark~\ref{rem:H3space}); the content of H3 is the
 exclusion of the non-vector-field generators, not a scalar test.  Stationarity on
 the linear class forces $(t_*Q)_{DD}=M_*$ self-adjoint (Prop.~\ref{prop:S}), i.e.\ S10's
 condition {\rm(S)} holds automatically at the new point.
\item \emph{Identification with S10} (Prop.~\ref{prop:s10}, Remark~\ref{rem:consist}):
 $\Sigma_2\to(tQ)_{DD}$, $\Sigma_1\to(tQ)_{DD}-t_{DD}$; $\Sigma_2\ge0$ \emph{is} the
 loss-of-isometry condition and $\Sigma_1\ge0$ \emph{is} the noise condition.
\item \emph{Modular-frame reduction of $\theta$} (Prop.~\ref{prop:thetaframe}): $\theta$ is
 a geometry-free Wiener--Hopf constant, with the explicit kernel \eqref{eq:kernelK} and
 $w(y)=-4\sinh^2(y/2)\mathbf 1_{y>0}$.  (Its \emph{value} is numerical input, not a
 theorem: see below.)
\end{enumerate}

\paragraph{Numerical input (not proved here).}
\begin{enumerate}[leftmargin=1.8em,itemsep=1pt]
\item $\theta=0.384\pm0.003$, $1-\theta=0.616\pm0.003$ (F3,
 \texttt{numerics/optimality\_all/F3\_RESULTS.md}, \texttt{f3\_t0.py}): both
 discretisations agree to $3\times10^{-4}$ after extrapolation in $h$ (order $h^{0.77}$);
 S10's $0.300$ is a taper-suppressed lower bound.  Remark~\ref{rem:twonum}.
\item Condition (a) holds in the limit, although every finite run fails it (F2,
 \texttt{numerics/\allowbreak optimality\_\allowbreak all/\allowbreak F2\_RESULTS.md}, T4): in the rigid frame of
 Prop.~\ref{prop:thetaframe}, $\lambda_{\min}(M_*)/\|M_*\|=-4.1\cdot10^{-3}$,
 $-1.9\cdot10^{-3}$, $-8.2\cdot10^{-4}$, $-3.7\cdot10^{-4}$ at $h=0.24,0.12,0.06,0.03$, the
 violation scaling as $h^{1.15}$ and as $1/Y_{\max}$, with rank-one gain $\to0$ as
 $h^{1.85}$ --- i.e.\ a discretisation artefact, exactly the configuration
 Cor.~\ref{cor:better} (R14) warns about, resolving \emph{in favour of} (a).
\item Condition (e) holds \emph{on the model cone} (F2 T5,
 \texttt{f2\_t5\_endpoint.out}; Remark~\ref{rem:T5}).  At $\zeta_\infty$:
 $\mathcal E(\zeta_\infty)$ reproduces $2\lambda(1-\theta_h)g_Q(\dot\delta)$ to $10^{-15}$
 (the proved half).  At $\zeta_0$: with the admissible (backward, dissipative) scheme the
 \emph{continuum piece} is the corner row
 $\frac12M_*(0,0)/h=+10.3,+16.1,+24.7,+37.2$ ($\sim h^{-0.60}$) --- the tabulated total
 $+34.1,+66.9,+131.3,+258.0$ ($\sim h^{-0.97}$) is dominated by the scheme-dependent upwind
 viscosity term of \eqref{eq:T5decomp} and must not be quoted as the margin.  Both diverge,
 so $\Aop(\zeta_0,0,0)\notin\mathfrak F$ and (e) is \emph{vacuous} at $\zeta_0$ rather than
 satisfied with a margin; the centred scheme gives $0$ (the linear class) and the forward
 (inadmissible) one the negative, i.e.\ the sign is decided by admissibility exactly as
 predicted.  No sign change in $h$, $Y_{\max}$ or the frame.  Nothing is known
 numerically about the non-vector-field generators of Remark~\ref{rem:H3}.
\end{enumerate}

\paragraph{Conditional.}
\begin{enumerate}[leftmargin=1.8em,itemsep=1pt]
\item \textbf{H1} ($g_Q$-differentiability along the scaled family,
 Def.~\ref{hyp:H1}): needed for the \emph{upper} bound $\limsup s^{-2}\min_\cF g_Q\le T$.
 Known for the isometric orbit (S10 Lemma~6.1) and for finite-rank $N_1,Y_1$ inside a
 modular window; open in general.  It fails, catastrophically, for the inadmissible
 contraction family of Remark~\ref{rem:coneT0} --- which is why admissibility is not a
 technicality.
\item \textbf{H2} (localisation of the minimisers, Def.~\ref{hyp:H2}): needed for the
 \emph{lower} bound, and it \emph{is} the lower bound in disguise (existence, scaled form,
 boundedness, uniformity).  Proposition~\ref{prop:coerc} gives $g_Q\ge\frac14\|\cdot\|_2^2$,
 hence $O(s)$ defect blocks, but Remark~\ref{rem:noloc} shows this does not localise
 $(X,Z)$: $Q_{AB}$ has $0$ in its essential spectrum and the invisible directions of
 Phase~4 are exactly the unresolved ones.  The unconditional substitute is the finite-$s$
 weak-duality certificate \eqref{eq:weakdual}.
\item \textbf{H3} (cone hypothesis, Def.~\ref{hyp:H3}, read in $\mathfrak F$): needed to
 pass from ``(e) on the model cone'' to ``(e)''.  Proved for vector-field generators; open
 for shift semigroups of multiplicity $>1$ and for unbounded skew-adjoint generators
 outside $\overline L$ (Remark~\ref{rem:H3}).  Its content is the exclusion of those two
 families: given it, (e) follows from $\mathcal E(\zeta_\infty)\ge0$ alone
 (Remark~\ref{rem:H3space}) if $\Aop(\zeta_0,0,0)\notin\mathfrak F$ (numerical, not proved; Correction~\ref{corr0924:C4}), the junction test \eqref{eq:econd} then being vacuous in
 $\mathfrak F$ and operative only for regularised directions.
\item Therefore: numerically (a) holds and (e) holds on the model cone, so \emph{under H3}
 the isometric orbit is the global optimum of the tangent problem and
 $\kappa_{\rm opt}=1-\theta=0.616\pm0.003$ --- modulo H1, H2, H3, and modulo the fact that
 (a) is extrapolated from meshes, on each of which it fails, and (e) is verified on meshes; neither is proved.  If (a) failed with a gain bounded away
 from $0$ (not what F2 sees), then $\kappa_{\rm opt}<1-\theta$ modulo H1 only.
\end{enumerate}

\paragraph{To be supplied.}
F2: delivered (T4 for (a), T5 for (e)).  Still useful: the cap- and grid-dependence
of $\sup_\psi\Delta^2/g_Q(d)$ \eqref{eq:gainreport}, to confirm that the residual violation
of (a) is an artefact; and $\|\mathrm{AntiHerm}((t_*Q)_{DD})\|$ as the M\"obius-rigidity
arbiter.  F3: done for $\theta$; what remains is a closed form, i.e.\ solving or bounding
the Wiener--Hopf problem of Prop.~\ref{prop:thetaframe} against $0.384\pm0.003$.  Analysis:
\textbf{H3} (evaluate $\mathcal E$ on multiplicity-$>1$ shift semigroups and on unbounded
skew-adjoint $\zeta=\mathrm im(y)$, or prove the cone reduction), H1 for general data, and
H2 (or the finite-$s$ certificate) for the lower bound.

\begin{verdictgap}[\;One-line status of the Phase-5 question]
The exact all-channel second-order value is $\kappa_{\rm opt}f_2z^2$, and the decision
whether $\kappa_{\rm opt}=1-\theta$ is no longer a search over channels but a pair of
computable form inequalities $M_*\ge0$, $M_*\ge\tau_*$ (condition (a); F2 finds them
satisfied in the continuum limit) \emph{together with} the endpoint condition (e) along the
admissible isometry generators (proved $\ge0$ at $\zeta_\infty$, strictly positive iff $\theta<1$, which is numerical only, Correction~\ref{corr0924:C2}; at $\zeta_0$ vacuous in
$\mathfrak F$, and positive and divergent along every admissible regularisation, F2 T5).
Both hold numerically \emph{on the model cone}, so with
$\theta=0.384\pm0.003$ (F3) the value is $\kappa_{\rm opt}=1-\theta=0.616\pm0.003$, modulo
H1, H2 and the cone hypothesis H3 (Def.~\ref{hyp:H3}), whose settlement --- or a direct
evaluation of $\mathcal E$ on multiplicity-$>1$ shifts and on unbounded skew-adjoint
generators --- is the sharpest remaining analytic task.
\end{verdictgap}

\section{Corrections of 2026-09-24}\label{corr0924:sec}
Each line of the text above that was changed was replaced by exactly one line, so every line
number of the version of 2026-09-23 still points to the same line.
%
\begin{enumerate}[label=C\arabic*,ref=C\arabic*,leftmargin=2.4em,itemsep=3pt]
\item\label{corr0924:C1} \textbf{Sign of $\mathcal G$ in Theorem~\ref{thm:bdry}.}\\
 Knowledge-base card \texttt{tangent-problem-exact-parametrisation}.
 \emph{Said:} $\mathcal G:=\zeta^*-\sigma$ in the first form of \eqref{eq:GS} (line 504), in
 Theorem~\ref{thm:bdry}(3) (line 520) and in step $\langle1\rangle4$ of its proof (line 552).
 \emph{Says now:} $\mathcal G:=\sigma-\zeta^*$ at all three places.
 \emph{Why:} the same display continues $\mathcal G=-\frac12(\zeta^*-\zeta^{**})$, and
 formally $\sigma=\frac12(\zeta^*+\zeta^{**})$, so $-\frac12(\zeta^*-\zeta^{**})=\sigma-\zeta^*$;
 the sentence after the display ($\zeta^*=-\mathcal G+\sigma$, $\zeta\subseteq\mathcal G+\sigma$)
 and step $\langle1\rangle2$ ($\zeta=\mathcal G+\sigma$, $\zeta^*=-\mathcal G+\sigma$) use this
 sign, and only with it does step $\langle1\rangle3$ give
 $\Aop_{AD}=-Q_{AD}\mathcal G+\frac12Q_{AD}(N_1-2\sigma)$ and
 $\Aop_{DD}=[\mathcal G,Q_{DD}]+\frac12\{N_1-2\sigma,Q_{DD}\}-Y_1$, i.e.\ the absorption
 identity \eqref{eq:absorb}.  With $\mathcal G=\zeta^*-\sigma$ the identity \eqref{eq:absorb}
 and the inclusion of Theorem~\ref{thm:bdry}(3) would be false as written.  Every other use of this $\mathcal G$ was
 checked and is unchanged: lines 506, 516, 539--556, 840, 969--984, 1205, 1533--1538, and
 Lemma~\ref{lem:deriv}, where $\mathcal G=-\lambda D_\chi|_D$ stands in the $\zeta$-slot of
 $\Aop$.  In particular ``$\mathcal G$ and $\zeta$ agree on $C_c^\infty(D)$'' (line 554) and
 ``$\sigma$ vanishes and $\mathcal G=\zeta$'' (line 984) hold only with the corrected sign,
 because $\sigma=0$ and $-\zeta^*=\zeta$ on $\operatorname{dom}\zeta$.  The letter
 $\mathcal G$ for the bounded anti-self-adjoint class (Cor.~\ref{cor:orbit}, with
 $\zeta=-\mathcal G$; Lemma~\ref{lem:commQ}; \S\ref{sec:f2}; \S\ref{sec:theta}) denotes a
 separately defined object and is not affected.
\item\label{corr0924:C2} \textbf{Strict positivity of the moving-endpoint term.}\\
 Knowledge-base card \texttt{kkt-noise-sign-criterion}.
 \emph{Said,} as proved: $\mathcal E(\zeta_\infty)=2\lambda(1-\theta)g_Q(\dot\delta)>0$,
 ``the moving-endpoint condition holds, unconditionally'' (Prop.~\ref{prop:bdryfun}(4),
 lines 1081--1082); the strict sign was repeated in Lemma~\ref{lem:deriv} (lines 718--721,
 ``It does not vanish \dots\ $>0$'', and proof step $\langle1\rangle6$, line 746, ``the
 counterexample''), after Theorem~\ref{thm:kktmain} (line 805, ``the moving-endpoint one is
 $>0$''), in step $\langle1\rangle5$ of the proof of Prop.~\ref{prop:bdryfun} (line 1131,
 ``the second term $>0$ for $\beta>0$''), in the Summary (line 1555, ``the moving-endpoint
 half is proved: \dots\ $>0$''; line 1558, ``only $\mathcal E(\zeta_\infty)>0$ is left'';
 line 1618, ``(e) follows from $\mathcal E(\zeta_\infty)>0$ alone'') and in the final
 verdict box (line 1643, ``proved positive at $\zeta_\infty$'').
 \emph{Says now:} $\mathcal E(\zeta_\infty)=2\lambda(1-\theta)g_Q(\dot\delta)\ge0$ at each
 of these places, with strict positivity equivalent to $\theta<1$ and marked as not proved;
 line 718 says that the vanishing of the derivative is not automatic outside $\overline L$
 (instead of ``It does not vanish''), and line 1127 adds the one-line reason for $\ge0$.
 \emph{Why:} step $\langle1\rangle4$ proves only the identity
 $\mathcal E(\zeta_\infty)=2\lambda g_Q(\delta_*)=2\lambda(1-\theta)g_Q(\dot\delta)$, hence
 $\mathcal E(\zeta_\infty)\ge0$ ($g_Q$ is a positive form, and $\theta\le1$ is trivial);
 $\mathcal E(\zeta_\infty)>0$ is equivalent to $\theta<1$, which this document records as
 \emph{not} proved (Remark~\ref{rem:whytheta}(i), lines 1461--1464; the verdict box after
 Remark~\ref{rem:twonum}, lines 1504--1507); numerically $1-\theta=0.616\pm0.003$
 \eqref{eq:thetaF3}.  No other conclusion changes, because the arguments that use this term
 need only $\ge0$: the model-cone reduction Prop.~\ref{prop:bdryfun}(5) (by additivity,
 $\mathcal E(\alpha\zeta_0+\beta\zeta_\infty)=\alpha\mathcal E(\zeta_0)+\beta\mathcal E(\zeta_\infty)$,
 so (e) on the model cone is $\mathcal E(\zeta_0)\ge0$ once $\mathcal E(\zeta_\infty)\ge0$),
 and (e) under Hypothesis~\ref{hyp:H3} (Remark~\ref{rem:H3space}:
 $\mathcal E(\zeta)=\beta\,\mathcal E(\zeta_\infty)\ge0$, as already written there and in
 Remark~\ref{rem:H3}).  The \emph{numerical} statements that $\mathcal E(\zeta_\infty)>0$ on
 every computed mesh (Remark~\ref{rem:T5}(iii); \S\ref{sec:summary}, Numerical input) are
 unchanged: they rest on the numerical $\theta_h<1$.  The same reading is recorded on
 knowledge-base card \texttt{kkt-noise-criterion-numerics} (under H3, (e) follows from the
 proved $\mathcal E(\zeta_\infty)\ge0$ alone).
\item\label{corr0924:C3} \textbf{Condition (e) is not proved on the model cone.}\\
 Knowledge-base cards \texttt{kkt-noise-sign-criterion} and \texttt{all-channel-optimum-value}.
 \emph{Said} (verdict box after Remark~\ref{rem:consist}, lines 955--957): ``(e), which is
 proved only on the model cone (Prop.~\ref{prop:bdryfun}) and reduces to
 $\mathcal E(\zeta_0)\ge0$ only under Hypothesis~\ref{hyp:H3}''.
 \emph{Says now:} (e) is not proved; on the model cone it is equivalent to
 $\mathcal E(\zeta_0)\ge0$ (Prop.~\ref{prop:bdryfun}(5), given $\mathcal E(\zeta_\infty)\ge0$
 by Prop.~\ref{prop:bdryfun}(4)), which holds only numerically, vacuously at $\zeta_0$
 (Remark~\ref{rem:T5}); beyond the model cone (e) needs Hypothesis~\ref{hyp:H3}, under which,
 read in $\mathfrak F$ and given $\Aop(\zeta_0,0,0)\notin\mathfrak F$ (numerical), (e) holds
 outright (Remark~\ref{rem:H3space}).
 \emph{Why:} Prop.~\ref{prop:bdryfun} proves the \emph{reduction} of (e) on the model cone,
 not (e) there: the junction inequality $\mathcal E(\zeta_0)\ge0$ is not proved, and F2's
 numerics only indicate $\mathcal E(\zeta_0)=+\infty$, the case
 $\Aop(\zeta_0,0,0)\notin\mathfrak F$ of \eqref{eq:Efun} (Prop.~\ref{prop:bdryfun}(5),
 Remark~\ref{rem:T5}(ii)).  Card \texttt{kkt-noise-sign-criterion} lists
 ``$\mathcal E(\zeta_0)>0$ on the model cone'' under ``numerical, not proved here'', and card
 \texttt{all-channel-optimum-value} lists an analytic proof of (e) as open.  The second half
 of the old sentence is restated in the cards' wording: the reduction to
 $\mathcal E(\zeta_0)\ge0$ holds on the model cone without H3; H3 is what extends (e) beyond
 the model cone, and under it the junction inequality is not invoked
 (Remark~\ref{rem:H3space}(1)).
\item\label{corr0924:C4} \textbf{$\Aop(\zeta_0,0,0)\notin\mathfrak F$ is numerical.}\\
 Knowledge-base card \texttt{kkt-noise-sign-criterion}.
 \emph{Said} (\S\ref{sec:summary}, ``Proved here, unconditionally'', item~5, lines
 1556--1557): ``Under Hypothesis~H3, read in $\mathfrak F$, (e) holds outright: since
 $\Aop(\zeta_0,0,0)\notin\mathfrak F$, \eqref{eq:H3} forces $\alpha=0$''.
 \emph{Says now:} (e) holds outright under H3 read in $\mathfrak F$ \emph{if}
 $\Aop(\zeta_0,0,0)\notin\mathfrak F$, which is numerical, not proved; then \eqref{eq:H3}
 forces $\alpha=0$.
 \emph{Why:} $\Aop(\zeta_0,0,0)\notin\mathfrak F$ rests on F2's divergent corner term
 (Remark~\ref{rem:T5}(ii)); Prop.~\ref{prop:bdryfun}(5) proves only that
 $-2\sigma_{\zeta_0}=\delta_0\otimes\delta_0$ is not in the form domain of $M_*$ and reports
 that the numerics indicate $\mathcal E(\zeta_0)=+\infty$.  The card states: ``Under H3, if
 $A(\zeta_0,0,0)$ is not in F then alpha = 0 is forced and (e) holds OUTRIGHT \dots; that
 $A(\zeta_0,0,0)$ is not in F \dots\ rests on F2's divergent corner term (rem:T5 (ii)), i.e.\
 it is numerical.''  The implication itself (H3 read in $\mathfrak F$ and
 $\Aop(\zeta_0,0,0)\notin\mathfrak F$ imply (e)) is proved by Remark~\ref{rem:H3space} and
 stays in the list.  The same proviso (``given $\Aop(\zeta_0,0,0)\notin\mathfrak F$,
 numerical'') was added, by one-line replacements, wherever the document invokes it for
 ``$\alpha=0$ is forced'' or ``(e) holds outright under H3'': line 1091
 (Prop.~\ref{prop:bdryfun}(5)); lines 1210 and 1213 (Remark~\ref{rem:H3space}; line 1213
 also says that Prop.~\ref{prop:bdryfun}(5) proves only that $\delta_0\otimes\delta_0$ is
 outside the form domain of $M_*$, the claim $\Aop(\zeta_0,0,0)\notin\mathfrak F$ resting on
 F2's divergent corner term, Remark~\ref{rem:T5}(ii)); lines 1232 and 1257
 (Remark~\ref{rem:H3}); line 1619 (\S\ref{sec:summary}, Conditional, H3).
\item\label{corr0924:C5} \textbf{The margin at $\zeta_\infty$ decreases under refinement.}\\
 Knowledge-base card \texttt{kkt-noise-criterion-numerics}.
 \emph{Said} (Remark~\ref{rem:T5}(iii), lines 1182--1183): (e) holds ``strictly and with an
 increasing margin at $\zeta_\infty$''.
 \emph{Says now:} strictly at $\zeta_\infty$, with a margin that decreases under refinement
 but stays positive, $\mathcal E(\zeta_\infty)/(\lambda g_Q(\dot\delta))=2(1-\theta_h)=1.37,
 1.32,1.28,1.26$ at $h=0.24,0.12,0.06,0.03$, towards $2(1-\theta)=1.232\pm0.006$.
 \emph{Why:} the card records, and
 \texttt{numerics/optimality\_all/f2\_t5\_endpoint.out} (lines 3, 13, 23, 33; $\lambda=1$)
 prints, $\mathcal E(\zeta_\infty)/g_Q(\dot\delta)=1.37265519$, $1.31527744$, $1.28132522$,
 $1.26109117$ at $h=0.24,0.12,0.06,0.03$, equal to $2(1-\theta_h)$ to $10^{-15}$ with
 $\theta_h=0.31367$, $0.34236$, $0.35934$, $0.36945$.  The margin falls as $h$ falls, because
 $\theta_h$ rises towards $\theta=0.384\pm0.003$ \eqref{eq:thetaF3}.  Only the word
 ``increasing'' was wrong; strictness and the absence of a sign change stand, as numerical
 statements.
\item\label{corr0924:C6} \textbf{The commutator form is not ``only solver noise''.}\\
 Knowledge-base card \texttt{kkt-noise-sign-criterion}.
 \emph{Said} (\S\ref{sec:f2}, item~5, line 1325): the commutator $-[Q,\zeta_0]$ ``returns
 only solver noise (Prop.~\ref{prop:bdryfun}(2$'$))''.
 \emph{Says now:} it carries no information about $\mathcal E$: it is not merely solver noise
 but $O(1)$ and worse on the finer scans.
 \emph{Why:} Prop.~\ref{prop:bdryfun}(2$'$), which the line cites, reports the F2 T5 values
 $-9.1\cdot10^{-6}$, $+1.9\cdot10^{-4}$, $-9.9\cdot10^{-3}$, $+7.4\cdot10^{-2}$ at
 $h=0.24,0.12,0.06,0.03$, other rungs reaching $-63$ and $-1.6\cdot10^6$, and says that the
 expression ``is not merely small noise, it is $O(1)$ and worse on the finer scans''; the card
 records the same values (``O(1), no information on the true $2\lambda(1-\theta)g_Q$'').
\item\label{corr0924:C7} \textbf{Agreement of the $\Aop$-form and the boundary form.}\\
 Knowledge-base card \texttt{kkt-noise-criterion-numerics}.
 \emph{Said} (Prop.~\ref{prop:bdryfun}(2$'$), line 1056, and Remark~\ref{rem:T5}(iii),
 line 1181): the two expressions agree to ``$2\cdot10^{-10}$--$2\cdot10^{-8}$''.
 \emph{Says now:} $2\cdot10^{-10}$--$7.5\cdot10^{-8}$.
 \emph{Why:} \texttt{numerics/optimality\_all/f2\_t5\_endpoint.out} prints, for the
 admissible (backward) scheme, the relative deviations $2.0\cdot10^{-10}$,
 $5.4\cdot10^{-9}$, $1.8\cdot10^{-8}$, $7.5\cdot10^{-8}$ at $h=0.24,0.12,0.06,0.03$ (lines 7,
 17, 27, 37), and values between $4.1\cdot10^{-9}$ and $6.5\cdot10^{-8}$ in the $Y_{\max}$-
 and $Y$-scans (lines 55--127); the card records ``holds to $\le7.5\cdot10^{-8}$
 (f2\_t5\_endpoint.out l.7, 17, 27, 37)''.
\item\label{corr0924:C8} \textbf{The corner row is not a margin.}\\
 Knowledge-base card \texttt{kkt-noise-criterion-numerics}.
 \emph{Said} (Remark~\ref{rem:T5}(ii), lines 1170--1171): the corner row ``is that row whose
 margin should be quoted''; (iii), line 1183: (e) holds ``vacuously/with a divergent positive
 corner value at $\zeta_0$''.
 \emph{Says now:} the corner row, not the total, is the row to quote for the regularised
 (finite-$h$) junction directions, as a divergent value and not as a margin; (e) holds at
 $\zeta_0$ vacuously, not with a margin (the corner value is positive at every $h$ and
 diverges).
 \emph{Why:} the corner row diverges ($\sim h^{-0.60}$), which the remark itself reads as
 $\Aop(\zeta_0,0,0)\notin\mathfrak F$, so that ``(e) is vacuous at $\zeta_0$, not that it
 holds there with a large margin'' (lines 1174--1175); by \eqref{eq:Efun} the value
 $+\infty$ is a convention, and reading it as a margin ``is legitimate only after (a) \dots\
 has been granted'' (lines 1033--1037), while (a) fails at every finite mesh.  The card says
 the corner row ``represents the continuum $E(\zeta_0)$'', ``diverges'', and that
 $E(\zeta_0)=+\infty$ is ``a convention, not a margin''.
\item\label{corr0924:C9} \textbf{Condition (a) is an extrapolation.}\\
 Knowledge-base card \texttt{kkt-noise-criterion-numerics}.
 \emph{Said} (\S\ref{sec:summary}, Numerical input, line 1577): ``Condition (a) holds for the
 exact problem''; (Conditional, line 1624): ``(a) and (e) are verified on meshes, not
 proved''.
 \emph{Says now:} ``Condition (a) holds in the limit, although every finite run fails it'';
 ``(a) is extrapolated from meshes, on each of which it fails, and (e) is verified on meshes;
 neither is proved''.
 \emph{Why:} the card: every finite run prints ``CRITERION (M$>$=0 and M-tau$>$=0): False''
 (f2\_t4\_h024.out l.4, f2\_t4\_hladder.out l.4 and l.12, f2\_t4\_h003.out l.4), and the
 violation only extrapolates to $0^-$ ($\lambda_{\min}(M_*)/\|M_*\|\sim h^{1.15}$ and
 $\sim1/Y_{\max}$): ``READ THE SIGN CLAIM AS AN EXTRAPOLATION''.  So (a) is not verified on
 any mesh. Line 1578 only gained \texttt{\textbackslash allowbreak} points inside its file name, because the longer line 1577 made that line overflow; the rendered text is unchanged.
\item\label{corr0924:C10} \textbf{Right-hand side of Theorem~\ref{thm:bdry}(3) in the
 abstract and in Remark~\ref{rem:coneT0}.}\\
 Knowledge-base card \texttt{tangent-problem-exact-parametrisation}.
 \emph{Said} (abstract, line 28): ``the cone generated by the anti-symmetric isometry
 generators''; (Remark~\ref{rem:coneT0}, line 967): ``Dropping `and an isometry
 generator'\,''.
 \emph{Says now:} ``the cone generated by the anti-symmetric parts of admissible isometry
 generators''; ``Dropping `the anti-symmetric part of an admissible $\zeta$' (i.e.\ admitting
 every anti-symmetric $\mathcal G$)''.
 \emph{Why:} both lines had the wording of Theorem~\ref{thm:bdry}(3) from before the third
 referee's point D2.  The amended theorem describes the right-hand side by how $\mathcal G$
 arises, as the anti-symmetric part of an admissible $\zeta$, and does not claim that
 $\mathcal G$ is itself an isometry generator (lines 519--525; proof step $\langle1\rangle4$).
 The card: ``G ... the anti-symmetric part of an ADMISSIBLE zeta, i.e.\ G is described by how
 it arises; that G itself generates an isometry semigroup is not claimed''.
\end{enumerate}
%
\paragraph{Remark.}  The places that invoke Remark~\ref{rem:H3space} for ``$\alpha=0$ is
forced'' or ``(e) holds outright under H3'' (lines 1091, 1210, 1213, 1232, 1257 and 1619) now
carry the proviso that $\Aop(\zeta_0,0,0)\notin\mathfrak F$ is numerical, not proved (C4).

\end{document}
